\documentclass[a4paper,11pt]{article}% {{{
% -------------------------------------------------
% {{{ Set up the page margins.
% -------------------------------------------------
\usepackage[margin=1in]{geometry}
% \usepackage[headheig\label{key}ht = .5in, headsep = \baselineskip, top = 1in, left = 1in, right = 1in, textwidth = 7.52 in]{geometry}
% \usepackage[text={6.5in,9.5in},centering]{geometry}
% \usepackage[margin=1in]{geometry}
% \setlength{\topmargin}{-0.25in}
% }}}
% -------------------------------------------------
% {{{ Load Packages here
% -------------------------------------------------
\usepackage[utf8]{inputenc}
\usepackage[T1]{fontenc}
\usepackage{amssymb,
  nccmath,
  geometry,
    latexsym,
    amsmath,
    graphics,
    fullpage,
    epsfig,
    amsthm,
    relsize,
    tikz,
    amsfonts,
    mathrsfs,
    makeidx,
    comment,
    latexsym,
    dcolumn,
    tocloft, % spacing for tableofcontents
    calc}
\usepackage[colorlinks,
     citecolor=blue,
     urlcolor=blue
      ]{hyperref}
\usepackage[shortlabels]{enumitem}
\excludecomment{codes}
% }}}
% -------------------------------------------------
% {{{ check references
% -------------------------------------------------
\usepackage{refcheck}
% \usepackage[notref,notcite]{showkeys}
% \usepackage[notcite]{showkeys}
% \usepackage{showkeys}
% }}}
% -------------------------------------------------
% {{{ Biber related
% -------------------------------------------------
\usepackage[strict=true,
            style=english]{csquotes}
% \usepackage{mathscinet}
\DeclareTextCommand{\polhk}{OT4}{\k}
\DeclareTextCommand{\polhk}{T1}{\k}

% \setlength{\parskip}{0.5cm plus4mm minus3mm}
% \usepackage[style=alphabetic, %authoryear
%             bibstyle=authoryear,
%             citestyle=authoryear,
%             natbib=true,
%             hyperref=true,
%             backref=true,
%             abbreviate=true]{biblatex}
\usepackage[backend=biber,
     style=alphabetic,
      hyperref=true,
      backref=true,
     natbib=true,
      sorting=nyt,
      abbreviate=true]{biblatex}
\AtEveryBibitem{% Clean up the bibtex rather than editing it
 \clearfield{doi}
 \clearfield{url}
 \clearfield{issn}
 \clearlist{location}
 \clearfield{month}
 \clearfield{series}

 \ifentrytype{book}{}{% Remove publisher and editor except for books
  \clearlist{publisher}
  \clearname{editor}
 }
}
% \addbibresource{All.bib}
\addbibresource{./draft_rand_int_biber.bib}
% }}}
% -------------------------------------------------
% {{{ Environments and some setups
% -------------------------------------------------
\theoremstyle{plain}
\newtheorem{theorem}{Theorem}[section]
\newtheorem{corollary}[theorem]{Corollary}
\newtheorem{lemma}[theorem]{Lemma}
\newtheorem{proposition}[theorem]{Proposition}

\theoremstyle{definition}
\newtheorem{definition}[theorem]{Definition}
\newtheorem{hypothesis}[theorem]{Hypothesis}
\newtheorem{assumptions}[theorem]{Assumption}

\newtheorem{remark}[theorem]{Remark}
\newtheorem{conjecture}[theorem]{Conjecture}
\newtheorem{example}[theorem]{Example}
\newtheorem{notation}[theorem]{Notation}

\numberwithin{equation}{section}
\allowdisplaybreaks
% }}}
% -------------------------------------------------
% {{{ Define short hand symbols.
% -------------------------------------------------
\newcommand{\lip}{\textup{Lip}_b}
\newcommand{\liph}{\text{Lip}_{\widehat{\rho} }}
\newcommand{\R}{\mathbb{R}}
\newcommand{\E}{\mathbb{E}}
\newcommand{\e}{\mathrm{e}}
\newcommand{\Norm}[1]{\left\|  #1   \right\|}
\newcommand{\rhoNorm}[1]{\left|  #1   \right|}
\newcommand{\ud}{\ensuremath{ \mathrm{d} }}
\newcommand*{\one}{ {{\rm 1\mkern-1.5mu}\!{\rm I} }}
% }}}
% }}}
% {{{ Head: Title, Abstract, Keywords
\title{Invariant measures for the stochastic heat equation on $\R^d$ without drift term}
\author{Le Chen\\Auburn University\\\texttt{le.chen@auburn.edu}\\
  \and
    Nicholas Eisenberg\\Auburn University\\\texttt{nze0019@auburn.edu}
  }
\date{\today}

\begin{document}
\maketitle

\begin{abstract}
  This paper deals with the long term behavior of the solution to the nonlinear
  stochastic heat equation $\partial u /\partial t - \frac{1}{2}\Delta u=
  b(u)\dot{W}$ for a Lipschitz continuous diffusion coefficient $b$ which is
  not necessarily bounded such as $b(u)=u$. Using the theory of the stochastic
  integral laid out by John Walsh, we provide conditions which will guarantee
  the existence of an invariant measure for a broad range of initial conditions
  which includes $L^2_\rho$ as well as the Dirac delta distribution $\delta_0$.
\end{abstract}
\bigskip

\noindent{\it \noindent MSC 2010 subject classification}: 60H15, 60H07, 60F05.
\smallskip

\noindent{\it Keywords}: Stochastic heat equation, parabolic Anderson model, invariant measure, Dirac delta initial condition.
\smallskip

% \noindent{\it Running head:} Ergodicity and CLT for PAM.

% \tableofcontents

% }}}
\section{Introduction}
We consider the following {\it stochastic heat equation} (SHE):
\begin{equation} \label{E:SHE}
    \begin{cases}
     \dfrac{\partial u}{\partial t}(t,x)-\dfrac{1}{2}\Delta u(t,x) = b(x,u(t,x))\dot W(t,x) & \text{$x\in \R^d$ , $t>0$}, \\
     u(0,\cdot) = \mu(\cdot),                                                               &
   \end{cases},
\end{equation}
The noise, $\dot{W}(t,s)$, is a centered Gaussian noise that is white in time
and homogeneously colored in space defined on a probability space $(\Omega,
\mathcal{F}, \mathbb{P})$ with natural filtration $(\mathcal{F}_t)$:
 \[
  \mathcal{F}_t = \sigma(W_s(A): \; 0 \le s \le t, \; A \in \mathscr{B}_b(\R^d)) \vee \mathcal{N},
 \]
where $\mathscr{B}_b(\R^d)$ denotes the bounded Borel sets on $\R^d$ and
$\mathcal{N}$ are the $\mathbb{P}$-null sets of $\Omega$. Its covariance
structure, $J$, is defined on $C_c^\infty(\R^{d+1}) \times
C_c^\infty(\R^{d+1})$ as follows:
\begin{align} \label{E:Cor}
  J(\psi,\phi) := \E \left[ \dot{W}\left(\psi\right)\dot{W}\left(\phi\right) \right]
                = \int_0^\infty \ud s \int_{\R^d} \Gamma(\ud x)(\psi(s,\cdot)*\widetilde{\phi}(s,\cdot))(x),
\end{align}
where $\widetilde{\phi}(x) := \phi(-x)$, ``$*$" refers to the convolution in
spatial variable, and $\Gamma$ is a nonnegative and nonnegative definite
tempered measure on $\R^d$ that is commonly referred to as the {\it correlation
measure}. The Fourier transform of $\Gamma$ (in the generalized sense) is also
a nonnegative and nonnegative definite tempered measure, which is called {\it
spectral measure} and is denoted by $\widehat{f}(\ud \xi)$. Moreover, in the
case where $\Gamma$ has density $f$, namely, $\Gamma(\ud x)=f(x) \ud x$, then
$\widehat{f}(\ud \xi) = \widehat{f}(\xi) \ud \xi$. The initial condition $\mu$
is a deterministic, locally finite, regular, signed Borel measure that
satisfies some integrability condition at the infinity (see \eqref{A:icon}
below). Let $\mu=\mu_+-\mu_-$ be its {\it Hahn decomposition} and
$|\mu|=\mu_++\mu_-$.

Any jointly measurable $L^2(\Omega)$-continuous $\mathcal{F}_t$-adapted
processes, $u=(u(t,x): \; t \ge 0,\; x\in \R^d)$, that is solves almost surely
the following stochastic integral equation is referred to as the {\it mild
solution} to \eqref{E:SHE}:
\begin{equation} \label{E:gmsol}
  u(t,x)=\int_{\R^d}G(t,x-y)\mu(\ud y) + \int_0^t\int_{\R^d}b(y,u(s,y))G(t-s,x-y)W(\ud s,\ud y),
\end{equation}
where $G(t,x) = (2 \pi t)^{-d/2}\exp\left( -\left(2t\right)^{-1}|x|^2\right)$
is the heat kernel and the stochastic integral is the {\it Walsh integral} (see
\cite{walsh:86:introduction,dalang.khoshnevisan.ea:09:minicourse}).

Regarding the function $b(t,x)$, we assume that for some constants $L_b>0$ and
$L_0\ge 0$, it holds that
\begin{equation} \label{E:lipcon}
  |b(x,u)-b(x,v)|< L_b |u-v| \quad \text{and} \quad
  |b(x,0)| \le L_0 \quad \text{for all } u,v \in \R, \; x \in \R^d.
\end{equation}
In particular, our assumption allows the linear case $b(x,u) = \lambda u$,
which is usually referred to the \textit{parabolic Anderson model}
\cite{carmona.molchanov:94:parabolic}. \bigskip

The aim of this paper is to investigate the conditions required to guarantee
the existence of invariant measures for the solution to \eqref{E:SHE}, which is
a crucial step towards the study of the ergodicity of the system. Let $\rho$ be
an admissible weight function (see Definition \ref{D:Rho}) and we will work
under the Hilbert space $H=L_\rho^2(\R^d)$. Let $\mathcal{B}(H)$ be the space
of all bounded Borel functions on $H$. Following
\cite{da-prato.g-atarek.ea:92:invariant}, a probability measure $\eta$ on the
Borel $\sigma$-field $\mathcal{B}(H)$ is said to be invariant for \eqref{E:SHE}
if
\begin{equation}\label{E:InvMeas}
  \eta(A) = \int_{H} \mathscr{L}(u(t,\cdot;\mu))(A) \; \eta(\ud  \zeta),
  \quad \text{for all $ t>0$ and $A\in \mathcal{B}(H)$},
\end{equation}
where we use the notation that $u(t,x;\varphi)$ with a third argument to
emphasize the dependence of the solution to \eqref{E:SHE} on the initial
condition $u(0,\cdot)=\varphi$. Note that in \eqref{E:InvMeas} we require $t$
to be strictly bigger than zero, unlike Eq. (3) of
\cite{da-prato.g-atarek.ea:92:invariant} where one imposes $t\ge 0$.


% We will see
% that for any $t_0>0$ that the solution of \eqref{E:SHE}, $u(t_0,\cdot)$,
% belongs to the Hilbert space $\mathcal{E}_{t_0} = (E_{t_0},
% \Norm{\cdot}_{E_{t_0}})$; see Remark \ref{R:Restart} below. Then by an
% invariant measure for \eqref{E:SHE}, we mean a measure, $\eta$, on
% $\mathscr{B}(E_{t_0})$ such that
% where $A \in \mathscr{B}(L^2_{\rho}(\R^d))$ and
% $\mathscr{L}(u(t,\cdot,\varphi))(A)$ is the law of \eqref{E:SHE} starting from
% $\varphi \in E_{t_0}$;
% \begin{equation} \label{D:law}
%  \mathscr{L}(u(t,\cdot,\varphi))(A) = \mathbb{P}\left[\omega \in  \Omega: u(t, \cdot, \varphi)(\omega) \in A\right], \quad t \ge t_0.
% \end{equation}
Due the {\it Krylov-Bogoliubov theorem} (see \cite[Theorem
3.1.1]{da-prato.zabczyk:96:ergodicity}), we see that that existence of the
measure can depend on the initial condition and under suitable conditions, for
example the assumptions of Theorem \ref{T:Main} below, the following measure,
$\eta$, is invariant for \eqref{E:SHE}:
\begin{equation} \label{E:InvMeasForm}
  \eta(A) = \lim_{n \to \infty} \frac{1}{T_n} \int_{t_0}^{T_n + t_0} \mathscr{L}(u(t,\cdot, \varphi))(A) \ud t, \quad \text{for any } t_0 >0.
\end{equation}
In the above, $\{T_n\}$ is an appropriately chosen sequence of positive numbers
increasing to infinity.

We note that proving uniqueness of the invariant measure is a much harder
question. An open question is if $\varphi_1$ and $\varphi_2$ are two such
initial conditions that are suitable for \eqref{E:InvMeasForm} and if $\eta_1$
and $\eta_2$ are the two invariant measures that arise from $\varphi_1$ and
$\varphi_2$ respectively then does $\eta_1 = \eta_2$? Another question one can
ask in the absence of uniqueness is what initial data are applicable to
\eqref{E:InvMeasForm} and can we give explicit examples? We address the second
question in Examples \ref{Ex:Gamma+f} and \ref{Ex:More} and also in the
discussion below.

% DaPrato and Zabczyk \cite{da-prato.g-atarek.ea:92:invariant} established the following procedure:
% \begin{enumerate}
%  \item show the solution is Markovian and that its transition semigroup is Feller,
%  \item show that the solution starting from a suitable initial condition, $\varphi \in E$, is bounded bounded in probability with respect to $E$.
% \end{enumerate}
% The first point does not restrict us when studying \eqref{E:SHE} as was shown in \textcolor{blue}{Inster Reference}. However the second point takes some work so show. As we will see below, in order to give actual examples, one needs to prove the stronger condition \eqref{E:BddMnt} in order to verify the boundedness in probability condition.

Existence in unbounded domains has been studied in few papers such as
\cite{assing.manthey:03:invariant, misiats.stanzhytskyi.ea:20:invariant,
misiats.stanzhytskyi.ea:16:existence,tessitore.zabczyk:98:invariant}. Since the
spatial domain is unbounded, one needs to introduce a weight function $\rho$
and the corresponding weighted space $L_\rho^2(\R^d)$ as usual; see Section
\ref{S:Prelim} below for more details. Following the standard procedure as laid
out by Da Prato and Zabczyk \cite{da-prato.g-atarek.ea:92:invariant}, the main
effort is to show that:
\begin{align} \label{E:BddMnt}
  \sup_{t>0} \E\left(\Norm{u(t,\cdot)}_{L^2_\rho(\R^d)}^2\right)
  <\infty.
\end{align}

However, this solution is usually {\it intermittent}, namely, the probability
moments to \eqref{E:SHE} have a certain exponential growth in $t$. In
\cite{assing.manthey:03:invariant,misiats.stanzhytskyi.ea:20:invariant,misiats.stanzhytskyi.ea:16:existence},
a drift term $f(x,u)$ is included in \eqref{E:SHE} in a crucial way to help to
cancel the otherwise exponential growing moments. The absence of such a drift
term in this paper makes the problem more challenging. In order to have moments
bounded in time, as required in \eqref{E:BddMnt}, one has to first identify the
exact conditions, under which the moments to \eqref{E:SHE} do not possess
exponential growth. This question has been answered in \cite[Theorem 1.3 and
Lemma 2.5]{chen.kim:19:nonlinear}, where necessary and sufficient conditions
are given. More precisely, for bounded second moments one has to have the
spatial dimension $d \ge 3$, and in addition, the spectral measure, $\hat{f}$,
and Lipschitz constant, $L_b$, (the Lyapunov exponent of $b$) in
\eqref{E:lipcon} must satisfy the following:
\begin{equation} \label{E:U0finite}
  \Upsilon(0) := \int_{\R^d} \frac{\hat{f}(\ud \xi)}{|\xi|^2} <\infty \quad \text{and} \quad
  64L^2_{b}   <  \frac{1}{4\Upsilon(0)};
\end{equation}
see the Assumption \eqref{A:dalcon0} below.

In the setting of equation \eqref{E:SHE}, Tessitore and Zabczyk
\cite{tessitore.zabczyk:98:invariant} prove existence of an invariant measure
for \eqref{E:SHE} in $L^2_{\rho}(\R^d)$ under the assumption that there exists
a $\varphi \in L^2_{\rho}(\R^d) \cap L^2_{\widetilde\rho}(\R^d)$ where $\rho
(\widetilde\rho)^{-1} \in L^1(\R^d)$ and the solution starting from $\varphi$
is bounded in probability in $L^2_{\widetilde \rho}(\R^d)$ and also that the
spectral density $\hat{f}$ satisfies
\begin{equation} \label{E:CorTZ}
  \widehat{f} \in L^p(\R^d) \quad \text{where } \frac{d-2}{d} < \frac{1}{p};
\end{equation}
see Hypothesis 2.1 ({\it ibid.}). However as was illustrated in \cite[Theorem
3.3]{tessitore.zabczyk:98:invariant}, in order to apply this theorem, the
following additional assumptions were needed: $d \ge 3$ and the Lipschitz
constant of $b$ must satisfy $L(b)^{-2} > \textbf{l}$, where
 \begin{equation} \label{E:TZspec}
  \widetilde \Gamma := \left| \sqrt{ \mathcal{F}(f) }  \right| * \left| \sqrt{\mathcal{F}(f)}  \right| \quad \text{and} \quad \textbf{l}:= \frac{\mathbf{\Gamma}(d/2 -1)}{4 \pi^{d/2}} \int_{\R^d} \widetilde{\Gamma}(\ud \zeta) |\zeta|^{2-d},
 \end{equation}
where $\mathbf{\Gamma}$ is the Gamma function. With these assumptions, they
were able to prove that \eqref{E:SHE} starting from the constant $1$ satisfies
\eqref{E:BddMnt} and thus is bounded in probability, verifying the existence of
an invariant measure. Moreover, the measure takes the form of
\eqref{E:InvMeasForm} above with $\varphi = 1$.

The additional assumptions given by Tessitore and Zabczyk parallel our
assumptions \eqref{E:U0finite} and \eqref{E:LipLb}. However, it is not at all
trivial to check which spectral densities, $\hat{f}(\ud \xi)$, satisfy
\eqref{E:TZspec} and no such examples were given in
\cite{tessitore.zabczyk:98:invariant}. For example, it may not be obvious from
\eqref{E:TZspec} that one can not take the spectral density of the noise to be
the Riesz kernel, $\hat{f}(\xi) = |\xi|^{d-\alpha}$, for any $0 < \alpha < d$
but this can be easily seen from our condition \eqref{E:U0finite}. The
condition that $d \ge 3$ forces us to only consider a colored noise as in
\eqref{E:Cor} and so the loss of the ability to apply the Riesz kernel is
unfortunate. However, as we will see in Example \ref{Ex:Gamma+f} below, we are
able to apply the Bessel kernel, $B_s(x)$, with $s < d$, which has exponential
decaying tails and power growth at the origin. So in absence of the Riesz
kernel, we still may apply our results to a kernel, namely the Bessel kernel,
that has origin blow up.

Moreover by applying a restart argument (see Remark \ref{R:Restart} and Example
\ref{Ex:Gamma+f} below), our method allows us to verify \eqref{E:BddMnt}, and
in return \eqref{E:InvMeasForm}, for more general initial conditions
$\mu(\cdot)$ which include:
\begin{align*}
  \text{Dirac delta measure }\delta_0(\ud x) \quad \text{and} \quad \varphi(x)\ud x \text{ with } \varphi \in C^{\infty}(\R^d).
\end{align*}
Thus, we can expand our choice of initial data to include singular
distributions, which are outside the space of $L^2_{\rho}(\R^d)$. We will
expand on this in the following remark.
%\begin{fleqn}
%\begin{align}
% (1) \; &\mu(\ud x) \text{ is a deterministic, locally finite, regular, signed Borel measure with the} \nonumber \\
% &\text{ corresponding {\it Hahn decomposition} $\mu=\mu_+-\mu_-$} \label{E:DetInit} \\
% (2) \; &\mu(\ud x) = \varphi(x) \ud x \text{ where $\varphi$ is an $\mathcal{F}_0$ measurable process in $L^2(\Omega, L^2_{\rho}(\R^d))$,} \nonumber\\
% & \text{namely  $\mathbb{E} \left( \Norm{\varphi(\cdot)}_{ L^2_{\rho}(\R^d)}\right) < \infty$.} \label{E:RandInit}
%\end{align}
%\end{fleqn}

\begin{remark}[Restarted SHE] \label{R:Restart}
  Suppose that $u(t_0,\cdot, \mu)$ is the solution to \eqref{E:SHE} starting
  from an initial condition, $\mu$, that satisfies the assumptions of Theorem
  \ref{T:Lrho}. Then \eqref{E:uInLrho} and Theorem \ref{T:Solve} imply that
  $u(t_0,\cdot, \mu)$ belongs to the Hilbert space, $\mathcal{E}_{t_0} :=
  (E_{t_0}, \; \Norm{\cdot}_{E_{t_0}})$, where $E_{t_0}$ is the space of random
  fields, $\{\varphi_{t_0}(x)\}_{\{x\in \R^d\}}$, that are $\mathcal{F}_{t_0}
  \times \mathscr{B}(\R^d)$ measurable with the following finite norm:
  \begin{align*}
    \Norm{\varphi_{t_0}}_{E_{t_0}}^2 := \mathbb{E} \left(\Norm{\varphi_{t_0}(\cdot)}_{L^2_{\rho}(\R^d)}^2\right) < \infty.
  \end{align*}
  % Thus, $u$ can be thought of as a solution to \eqref{E:SHE} starting from a
  % process $\varphi_{t_0} \in E_{t_0}$.
  %
  Then, for any $t_0 >0$, consider the following restarted equation:
   \begin{equation} \label{E:restart}
    \begin{cases}
    \dfrac{\partial v}{\partial t}(t,x)-\dfrac{1}{2}\Delta v(t,x) = b(x,v(t,x))\dot W_{t_0}(t,x) & \text{$x\in \R^d$ , $t>0$}, \\
    v(0,\cdot) = u(t_0,\cdot),                                                                   &
    \end{cases}
   \end{equation}
  where $\dot{W}_{t_0}(t,x) := \dot{W}(t + t_0,x)$ denotes the time shifted
  noise, ie,
   \[
    \int_0^t \int_{\R^d} W_{t_0}(\ud s, \ud y) = \int_{t_0}^{t+t_0} \int_{\R^d} W(\ud s, \ud y).
   \]
  Moreover, the Markov property of the solution to \eqref{E:SHE} implies that
  for any $t\ge 0$, we have $u(t + t_0,x) \stackrel{\mathcal{L}}{=} v(t,x)$
  where $\mathcal{L}$ denotes equality in law. Since we are interested in
  studying the long time behavior of \eqref{E:SHE}, we obtain the same
  information by choosing to study the restarted equation.

  This is a very import observation. Recall that Tessitore et.al. \cite[Theorem 3.3]{tessitore.zabczyk:98:invariant} proved that there exists in invariant measure, $\eta$, on $\mathscr{B}(L^2_{\rho}(\R^d))$ of the form \eqref{E:InvMeasForm} for the constant initial data $\varphi = 1$. Because of the restart property of the equation, we may rewrite \eqref{E:InvMeasForm} as the following:
   \begin{align*}
    \eta(A) &= \lim_{n \to \infty} \int_{t_0}^{T_n + t_0} \mathscr{L}(u(t, \cdot, 1))(A) \;\ud t \\
    & = \lim_{n \to \infty} \int_{t_0}^{T_n + t_0} \mathscr{L}\left[v\left(t - \frac{t_0}{2}, \cdot, u\left(\frac{t_0}{2}, \cdot, 1 \right)\right)\right](A) \; \ud t \\
    &= \lim_{n \to \infty} \int_{t'_0}^{T_n + t'_0} \mathscr{L}\left[v\left(t, \cdot, u\left(\frac{t_0}{2}, \cdot, 1 \right)\right)\right](A) \; \ud t, \quad t_0' := t_0 /2,
   \end{align*}
  where we recall that $A \in \mathscr{B}(L^2_{\rho}(\R^d))$ and $v$ denotes the solution to \eqref{E:restart} starting from $u(t_0/2, \cdot, 1)$ with the time-shifted noise $\dot{W}_{t_0/2}$. Naturally the invariant measure is usually defined on the Hilbert space of functions where the initial condition lives; in the case of \cite{tessitore.zabczyk:98:invariant}, this is $L^2_\rho(\R^d)$. The above shows that we can extend the measure $\eta$ to $\mathscr{B}(E_{t_0})$ since we can associate the equation \eqref{E:SHE} starting from $\mu = 1$ with the above restarted process $v$ that starts from $u(t_0, \cdot, 1) \in E_{t_0}$. This is no contradiction since clearly $\mathscr{B}(L^2_{\rho}(\R^d)) \subset \mathscr{B}(E_{t_0})$.

  This observation will allow us to show existence of an invariant measure in $\mathscr{B}(E_{t_0})$ for our equation \eqref{E:SHE} starting from an initial condition outside of $L^2_{\rho}(\R^d)$. For example, we will see in Example \ref{Ex:Gamma+f} below that we may start our equation at the Dirac delta measure $\delta_0$, which is clearly not in $L^2_{\rho}(\R^d)$.  However, because of the smoothing properties of the heat operator, after any amount of time $t_0>0$, the solution $u(t_0, \cdot, \delta_0) \in E_{t_0}$. Thus, whether we originally start our equation \eqref{E:SHE} at the constant 1, as in \cite{tessitore.zabczyk:98:invariant}, or at a distribution that satisfies the assumptions of Theorem \ref{T:Main} main below, like the Dirac delta distribution $\delta_0$, then either way we end up in the same function space. If one were able to prove uniqueness of the invariant measure, then the  invariant measure which is derived from the Dirac delta initial condition and the invariant measure which is derived from the constant $1$ initial condition would be the same.

\end{remark}

%According to \cite[Theorem 1.3]{chen.kim:19:nonlinear}, in
%order to have \eqref{E:BddMnt}, one has to have $L_b$ (the Lyapunov exponent of $b$) in \eqref{E:lipcon}  small enough, the spatial dimension $d \ge 3$, and we must strengthen \eqref{A:dalcon} to the following:
%\begin{equation} \label{E:U0finite}
%\Upsilon(0) <\infty.
%\end{equation}

%\begin{assumptions} \label{A:General}
%\begin{enumerate}[(i)]
%    \item The initial measure $\mu$ must satisfy the following integrability condition:
%    \begin{equation} \label{A:icon}
%    \int_{\R^d}\exp\left(-a|x|^2\right)\Norm{\mu(\ud x)}_2 < \infty \quad \text{for all $a>0$}.
%    \end{equation}
%    \item The spectral measure $\widehat{f}$ satisfies {\it Dalang's condition}, namely,
%      \begin{equation} \label{A:dalcon}
% \Upsilon(\beta) := (2\pi)^{-d}\int_{\R^d}\dfrac{\widehat{f}(\ud \xi)}{\beta + |\xi|^2} < \infty
% \quad \text{for some and therefore any $\beta>0$}.
%      \end{equation}
%  \end{enumerate}
%\end{assumptions}

%\begin{remark} \label{R:Cor}
%The assumption on the correlation function in
%\cite{tessitore.zabczyk:98:invariant} is given through the spectral density
%$\widehat{f}$ as follows:
%  \begin{equation} \label{E:CorTZ}
%    \widehat{f} \in L^p(\R^d) \quad \text{where } \frac{d-2}{d} < \frac{1}{p};
%  \end{equation}
% see Hypothesis 2.1 ({\it ibid.}). However as was illustrated in \cite[Theorem
% 3.3]{tessitore.zabczyk:98:invariant}, \eqref{E:CorTZ} was not enough of an
% assumption to give a concrete example of the existence of an invariant
% measure and additional assumptions were needed. For example, one needs $d \ge
% 3$ and the Lipschitz constant of $b$ to satisfy $l^{-2} > \textbf{l}$, where
%   \begin{equation} \label{E:TZspec}
%    \widetilde \Gamma := \left| \sqrt{ \mathcal{F}(f) }  \right| * \left| \sqrt{\mathcal{F}(f)}  \right| \quad \text{and} \quad \textbf{l}:= \frac{\mathbf{\Gamma}(d/2 -1)}{4 \pi^{d/2}} \int_{\R^d} \widetilde{\Gamma}(\ud \zeta) |\zeta|^{2-d},
%   \end{equation}
% where $\mathbf{\Gamma}$ is the Gamma function. These additional assumptions
% parallel our assumptions \eqref{E:U0finite} and \eqref{E:LipLb}, however one
% can see that both \eqref{E:U0finite} and \eqref{E:LipLb} are more compact and
% easier to check than \eqref{E:TZspec}.
%\end{remark}

For the remainder of the paper, $\mu$ will always denote the initial condition
to \eqref{E:SHE} and $u(t,\cdot) = u(t,\cdot, \mu)$ will always denote the
solution of \eqref{E:SHE} starting from $\mu$. In addition, $v_{t_0}(t,\cdot) =
v(t,\cdot, u(t_0,\cdot, \mu))$ will always denote the solution to the restarted
equation \eqref{E:restart} starting from $u(t_0,\cdot, \mu)$ and driven by the
time shifted noise $\dot{W}_{t_0}$. When there is no confusion, we may write
$v$ instead of $v_{t_0}$.

\begin{notation} \label{N:StartNotation}
  Denote $\Norm{\cdot}_\rho := \Norm{\cdot}_{L^2_\rho(\R^d)}$ and $\Norm{X}_p
  := \left(\mathbb{E}(|X|^p)\right)^{1/p}$. In addition, we define
  \begin{gather} \label{E:Ju}
     J_u(t,x,\mu) := (G(t,\cdot) * \mu)(x) = \int_{\R^d} G(t,x-y) \mu(\ud y), \quad \text{and}\\
     J_u^*(t,x) := J_0(t, x, \mu^*)(t,x) \quad \text{with} \quad \mu^*:= 1+ |\mu|.
     \label{E:J*u}
  \end{gather}
\end{notation}

We first present a preliminary result which is used to prove Theorem
\ref{T:Main} below.

\begin{theorem} \label{T:Lrho}
  Let $u(t,x)$ be the solution to \eqref{E:SHE} starting from $\mu$. Assume
  that
  \begin{enumerate}[(i)]
    \item $\rho: \R^d \to \R_+$ is an nonnegative $L^1(\R^d)$ function;
    \item the initial measure $\mu$ satisfies
      \begin{equation} \label{E:G}
        \mathcal{G}^u_{\rho}(t)
        := \int_{\R^d} \left(J_u^*(t,x)\right) ^2 \rho(x)\: \ud x< \infty
        \quad \text{for all $t>0$};
      \end{equation}
    \item the correlation measure $\Gamma$ and the Lipschitz constant $L_b$ of
      $b$ satisfy conditions in \eqref{E:U0finite}.
  \end{enumerate}
  Then for all $t>0$, $u(t,\cdot)\in L_{\rho}^2(\R^d)$ a.s. Moreover, for some
  constant $C>0$ which does not depend on $t$, it holds that
  \begin{equation} \label{E:uInLrho}
    \E\left(\Norm{u(t,\cdot)}_{\rho}^2\right)
    \le C \mathcal{G}^u_{\rho}(t)
    < \infty.
  \end{equation}
\end{theorem}
\begin{proof}
  Under condition (iii), we can apply Fubini's Theorem and the moment bound
  \eqref{E:mb0} below to see that for some constant $C>0$ independent of $t$,
  which may vary from line to line, it holds that
  \begin{align*}
    \E\left(\Norm{u(t,\cdot)}_{\rho}^2\right)
    & = \E \left(\int_{\R^d}\vert u(t,x) \vert^2 \rho(x)\ud x\right)                \\
    & \le C \int_{\R^d}\bigg(1+(G(t,\cdot) * |\mu|)(x)\bigg)^2\rho(x)\ud x          \\
    & = C \int_{\R^d}\bigg( \Big(G(t,\cdot) *(1+|\mu|)\Big)(x) \bigg)^2\rho(x)\ud x \\
    & = C \:\mathcal{G}_{\rho}(t) < \infty.
  \end{align*}
  This proves Theorem \ref{T:Lrho}.
\end{proof}

Here is the main theorem of the paper:

\begin{theorem} \label{T:Main}
  Let $u(t,x)$ be the solution to \eqref{E:SHE} starting from $\mu$. Assume
  that all conditions in Theorem \ref{T:Lrho} are satisfied. Moreover, we
  assume that
  \begin{enumerate}[(i)]
    \item if $\rho$ is an admissible weight (see Definition \ref{D:Rho} below)
      and there exists another admissible weight $\widehat{\rho}$ such that
      \begin{align} \label{E:rhorho}
        \int_{\R^d}\dfrac{\rho(x)}{\widehat{\rho}(x)}\ud x < \infty;
      \end{align}
    \item the weight function $\widehat{\rho}$ and the initial condition
      satisfy the following condition:
      \begin{equation} \label{E:InitData}
        \limsup_{t>0}\: \mathcal{G}^u_{\widehat{\rho}}(t+ t_0) < \infty;
      \end{equation}
    \item for some $\alpha\in (0,1/2)$, the spectral measure $\widehat{f}$
      satisfies that
      \begin{align} \label{E:H}
         \mathcal{H}_\alpha(t):= \int_0^t \ud r \: r^{-2\alpha} \int_{\R^{d}} \widehat{f}(\ud \xi) \exp\left(-r|\xi|^2 \right) < \infty
         \quad \text{for all $t>0$};
      \end{align}
    \item for the constant $\alpha$ in (ii), there exists $q>1/\alpha$ (i.e.,
      $q^{-1}< \alpha< 2^{-1}$) such that
      \begin{align} \label{E:GH-int}
      \int_{0}^t \left[\mathcal{G}^u_{\widehat{\rho}}(s + t_0)\:\mathcal{H}_\alpha(s)\right]^{q/2} \ud s <\infty \quad \text{for all $t>0$};
      \end{align}
    \item the second condition in \eqref{E:U0finite} is strengthened
      accordingly with $q$ given in (iv) to
      \begin{align*}
        qL_b^2 < \frac{1}{128 \Upsilon(0)}.
      \end{align*}
  \end{enumerate}
  Then we have that
  \begin{enumerate}[(a)]
    \item for any $ t_0>0$, the sequence of laws of
      $\{\mathcal{L}u(t,\cdot;\mu)\}_{t\ge t_0}$ is tight, i.e., for any
      $\epsilon \in (0,1)$, there exists a compact set $\mathscr{K}\subset
      L^2_{\rho}(\R^d)$ such that
      \begin{equation} \label{E:cpt}
        \mathcal{L}u(t,\cdot;\mu)(\mathcal{K})
        :=  \mathbb{P}\left( u(t,\cdot;\mu) \in \mathscr{K}\: \right)
        \ge 1-\epsilon, \qquad \text{for all $t\ge t_0>0$};
      \end{equation}
    \item there exists an invariant measure for the laws
      $\{\mathcal{L}u(t,\cdot;\mu)\}_{t>0}$.
  \end{enumerate}
\end{theorem}

We will postpone the proof of Theorem \ref{T:Main} (a) to Section
\ref{S:mainproof}. The proof of part (b) based on (a) is routine:

\begin{proof}[Proof of Theorem \ref{T:Main} (b)]
  Fix an arbitrary $t_0>0$ and denote
  \begin{align*}
    U(T) := \dfrac{1}{T}\int_{t_0}^{T+t_0} \mathcal{L}\left(u(t,\cdot;\mu)\right)\: \ud t, \quad T>0.
  \end{align*}
  We claim that the family of laws $U(T,\cdot)$ for $T > 0$ is tight in
  $L^2_{\rho}(\R^d)$. Indeed, for any $\epsilon\in (0,1)$, by part (a), there
  exists a compact set $\mathcal{K}\in L_\rho^2(\R^d)$ such that \eqref{E:cpt}
  holds. This implies that
  \begin{align*}
    U(T)\left(\mathcal{K}\right)
    & = \frac{1}{T} \int_{t_0}^{T + t_0} \mathscr{L}(u(t,\cdot;\mu))\left(\mathcal{K}\right) \; \ud t
      \ge \frac{1}{T} \int_{t_0}^{T + t_0} (1-\epsilon) \; \ud t
      = 1-\epsilon, \quad \text{for all $T>0$}.
  \end{align*}

  Let $\{T_n\}_{n\in \mathbb{N}}$ be any deterministic sequence such that
  $T_n\uparrow \infty$. Since $\left\{U(T_n)\right\}_{n\ge 1}$ is a tight
  sequence of measures, there exists a subsequence
  $\{U\left(T_{n_m}\right)\}_{m\ge 1}$ that converges weakly to a measure, say,
  $\eta$, on $L_\rho^2(\R^d)$. Then one can apply for the Krylov-Bogoliubov
  theorem (see, e.g., \cite[Theorem 3.1.1]{da-prato.zabczyk:96:ergodicity}) to
  conclude that the measure $\eta$ is the invariant measure for
  $\mathscr{L}(u(t,\cdot;\mu))$, $t\ge t_0$. Finally, since $t_0$ can be
  arbitrarily close to zero, one can conclude part (b) of Theorem \ref{T:Main}.
\end{proof}


\section{Examples} \label{S:Example}

\begin{example}[Bessel kernel and $\delta_0$ initial data] \label{Ex:Gamma+f}
  We end this section with an example of an applicable initial data and spectral
  density. Note that because of \eqref{E:GH-int}, the initial data and spectral
  density must exist in tandem. Let's consider the initial data $\mu(\cdot) =
  \delta_0(\cdot)$ and let the correlation function of $W$ be $f(x) = B_s(x)$,
  where $B_s$ denotes the Bessel kernel with parameter $s$. It is shown in
  \cite[Section 1.2.2]{grafakos:14:modern} that there exists $c(s,d)$ and
  $C(s,d)$ such that for $0 < s < d$
   \[
    B_s(x) \le C(s,d) \exp(-|x|/2) \quad \text{for} \quad |x| \ge 2,
   \]
  and for $H(x) = |x|^{s-d} + 1 + O(|x|^{s-d+2})$,
   \[
    \frac{1}{c(s,d)} \le \frac{B_s(x)}{H_s(x)} \le c(s,d) \quad \text{for} \quad |x| \le 2.
   \]
  Hence $B_s(x)$ has blow up at the origin and exponential decay at infinity.
  Moreover,
  \begin{equation} \label{E:FB}
    \mathcal{F}G_s(\xi) = \frac{1}{(1+  |\xi|^2)^{s/2}}.
  \end{equation}
  It is easy to see that $\delta_0$ satisfies \eqref{A:icon}. Moreover, it is
  easy to see that if $d-2 < s < d$ then $\Upsilon(0) < \infty$ and so the
  conditions of Theorem \eqref{T:Solve} and Theorem \ref{T:ExUn} are satisfied
  and a unique mild solution, $u$, exists that starts from $\delta_0$ and is
  driven by a Gaussian noise with correlation function $B_s(x)$. We denote
  $v(t,x)$ to be the restarted solution as in \eqref{E:restart}. We now show that
  $v$ satisfies the conditions in Theorem \ref{T:Main}. In other words, it
  remains to prove that $v(0,\cdot) = u(t_0,\cdot)$ and $B_s(x)$ satisfy the
  conditions \eqref{E:InitData}, \eqref{E:H} and \eqref{E:GH-int}.

  Let $q>2$ be arbitrary. Starting with \eqref{E:InitData},
   \begin{align*}
    \mathcal{G}_{\widehat{\rho}, q}(s)
    & =  \int_{\R^d} \left( \int_{\R^d} G(s,x-y) \mu_q^*(\ud y)\right)^2 \rho(x) \ud x                        \\
    & = \int_{\R^d} \left( \int_{\R^d} G(s,x-y) (1 + \Norm{v(0,y)}_q)(\ud y)\right)^2 \widehat{\rho}(x) \ud x \\
    & = \int_{\R^d} \left( 1 + \int_{\R^d} G(s,x-y) \Norm{v(0,y)}_q\ud y\right)^2 \widehat{\rho}(x) \ud x     \\
    & \le 2 \Norm{\widehat{\rho}}_{L^1(\R^d)} + 2\int_{\R^d} \left(\int_{\R^d} G(s,x-y) \Norm{u(t_0,y)}_q\ud y\right)^2 \widehat{\rho}(x) \ud x,
   \end{align*}
  where in the last line we used Recall that $u(t_0,x) = v(0,x)$. In the
  following, $C$ is a generic constant independent of $s$ that may change from
  line to line. By applying the moment bound \eqref{E:mb0} we see that
   \begin{align*}
     \mathcal{G}_{\widehat{\rho}, q}(s)
     & \le 2 \Norm{\widehat{\rho}}_{L^1(\R^d)} + 2C \int_{\R^d} \left(\int_{\R^d} G(s,x-y) \left[1 + (G(t_0,\cdot)*\delta_0)(y)\right]\ud y\right)^2 \widehat{\rho}(x) \ud x \\
     & = 2 \Norm{\widehat{\rho}}_{L^1(\R^d)} + C \int_{\R^d} \left(1  +\int_{\R^d} G(s,x-y) (G(t_0,\cdot)*\delta_0)(y) \ud y\right)^2 \widehat{\rho}(x) \ud x                \\
     & =2 \Norm{\widehat{\rho}}_{L^1(\R^d)} + C \int_{\R^d} \left(1  + (G(t_0+s,\cdot)*\delta_0)(x) \right)^2 \widehat{\rho}(x) \ud x                                        \\
     & \le 2 \Norm{\widehat{\rho}}_{L^1(\R^d)} + 2C\Norm{\widehat{\rho}}_{L^1(\R^d)} +  2C\int_{\R^d} (G(t_0+s,\cdot)*\delta_0)(x)^2 \widehat{\rho}(x) \ud x                 \\
     & \le C\Norm{\widehat{\rho}}_{L^1(\R^d)} +  C\int_{\R^d} (G(t_0+s, x))^2 \widehat{\rho}(x) \ud x                                                                        \\
     & = C\Norm{\widehat{\rho}}_{L^1(\R^d)} + C\frac{(2\sqrt{\pi})^{-d}}{(t_0 + s)^{d/2}} \int_{\R^d} G\left(\frac{t_0 + s}{2}, x\right) \widehat{\rho}(x) \ud x             \\
     & =: \textbf{G}(s),
   \end{align*}
  and this verifies \eqref{E:InitData} because $t_0 > 0$ is fixed. Moreover, this
  implies that for any $T>0$, $\mathcal{G}_{\widehat{\rho}, q}(s)$ is bounded on
  $s \in [0,T]$. This will be needed below.

  Now lets verify \eqref{E:H} for the Bessel kernel, $B_s$. By recalling
  \eqref{E:FB}, we have
   \begin{align*}
    \mathcal{H}_{\alpha}(s)
    & = \int_0^s \ud r \; r^{-2\alpha} \int_{\R^d} \ud \xi \; \frac{\exp(-r|\xi|^2)}{ (1+|\xi|^2)^{s/2}} \\
    & = \sigma\left(\mathbb{S}^{d-1}\right)\int_0^t \ud r \; r^{-2\alpha} \int_0^\infty \ud z \; \frac{\exp(-r z^2)}{(1+z^2)^{s/2}} z^{d-1}.
   \end{align*}
  Let's examine the $\int_0^\infty \ud z$ integral and apply the change of
  variables $y = z^2$. This yields,
  \begin{align*}
      \int_0^\infty \ud z \; \frac{\exp(-r z^2)}{(1+z^2)^{s/2}} z^{d-1}
    = \frac{1}{2} \int_0^\infty \ud y\; e^{-r y} (1+y)^{s/2} y^{\frac{d}{2}-1} \ud y
    = \frac{\Gamma\left(d/2\right)}{2}U\left(\frac{d}{2},\frac{d+2-s}{2},r\right),
  \end{align*}
  where $U\left(\alpha,\beta,t\right)$ is the confluent hypergeometric U function
  (see \cite[Chapter 13]{olver.lozier.ea:10:nist}) and the last equality can be
  found in \cite[13.4.4]{olver.lozier.ea:10:nist}. Thus, to show the finiteness
  of $\mathcal{H}_{\alpha}(s)$ for each $s >0$, it suffices to show that the
  following integral is finite:
   \begin{equation} \label{E:BesselHG}
    \int_0^s r^{-2\alpha} \; U\left(\frac{d}{2},\frac{d+2-s}{2},r\right) \; \ud r.
   \end{equation}
  The behavior of $U(\alpha, \beta, t)$ around the origin is given in
  \cite[Section 13.2 (iii)]{olver.lozier.ea:10:nist} and is given in cases on
  $\beta$, which in our case is $(d+2-s)/2$. Recall above that we are taking $d-2
  < s < d$. For this reason, the only case in \cite[Section 13.2
  (iii)]{olver.lozier.ea:10:nist} that is applicable is $1 < (d+2-s)/2 < 2$ and
  in this case
  \begin{equation}
    r^{-2\alpha} U\left(\frac{d}{2},\frac{d+2-s}{2},r\right) = O\left( r^{1-2\alpha- (d+2-s)/2} + r^{-2\alpha} + r^{2 - 2\alpha -  (d+2-s)/2}  \right).
  \end{equation}
  So in order \eqref{E:BesselHG} to be finite we need
   \[
    \min\{ 1 - 2\alpha -  (d+2-s)/2, -2\alpha, 2 - 2\alpha -  (d+2-s)/2\} \ge -1.
   \]
  Recall that we are assuming $\alpha \in (q^{-1}, 2^{-1})$ so $-2\alpha >-1$.
  Thus it suffices to set
   \[
    1 - 2\alpha -  (d+2-s)/2, -2\alpha \ge -1 \quad \iff \quad 2 - 4\alpha -d + s >0.
   \]
  This is indeed possible as $2-4\alpha >0$ and $d-2 < s < d$, thus one can take
  $s$ sufficiently close to $d$ so that $(2 - 4\alpha) - (d-s) >0$. This
  completes the verification for \eqref{E:H}.

  Lastly we need to verify \eqref{E:GH-int}. But essentially all the work for
  this has been completed. The verification of \eqref{E:H} implies that
  $\mathcal{H}_\alpha(s)$ is continuous on $s \in [0, T]$ for any $T>0$ and we
  already saw aboves that $\mathcal{G}_{\widehat{\rho}, q}(s)$ was bounded on
  $[0, T]$. These two facts imply \eqref{E:GH-int}, thus allowing us to apply
  Theorem \eqref{T:Main} and verify the existence of an invariant measure to
  \eqref{E:SHE} with initial data, $\mu(\cdot) = u(t_0, \cdot)$ and noise
  $\dot{W}_{t_0}$ with correlation function given to be the Bessel kernel $B_s$.
  Recall that the invariant measures derived from specific intial conditions take
  the form of \eqref{E:InvMeasForm}. Since for all $t >0$ we have $u(t + t_0,x)
  \stackrel{\mathcal{L}}{=} v(t,x)$, what we have actually done was verify the
  existence of an invariant measure$, \eta$, for \eqref{E:SHE} starting from
  $\delta_0$ with correlation function $B_s(x)$. Indeed,
   \begin{align*}
    \eta(A) & = \lim_{n \to \infty} \int_{t_0}^{T_n + t_0} \mathscr{L}(v(t,\cdot, u(t_0,\cdot)))(A) \; \ud t   \\
            & = \lim_{n \to \infty} \int_{t_0}^{T_n + t_0} \mathscr{L}(u(t + t_0,\cdot, \delta_0))(A) \; \ud t \\
            & = \lim_{n \to \infty} \int_{t'_0}^{T_n + t'_0} \mathscr{L}(u(t,\cdot, \delta_0))(A) \; \ud t, \quad t'_0 = 2t_0.
   \end{align*}
  Hence we may start \eqref{E:SHE} at $\delta_0$.
\end{example}

\begin{example}[More examples] \label{Ex:More}
We would like to point out that the following initial data are also applicable
when considering the correlation function to be the Bessel kernel $B_s(x)$ with
$s < d$:
 \[
  \text{Lebesgue measure } \ud x \quad \text{and} \quad \varphi(x)\ud x \text{ with } \varphi \in C^{\infty}(\R^d).
 \]
The interested reader may go trough the same steps above to see that
\eqref{E:InitData}, \eqref{E:H} and \eqref{E:GH-int} are all finite.
\end{example}

\begin{example}[Bounded initial conditions]
  When the initial condition is bounded, namely, $\mu(\ud x) =\phi(x)\ud x$
  with $\phi\in L^\infty(\R^d)$, our conditions in our Theorems \ref{T:Lrho}
  and \ref{T:Main} will be greatly simplified. \textcolor{magenta}{(Give
  details here...)}
\end{example}
\section{Preliminaries} \label{S:Prelim}

% \subsection{The weighted space $L^2_\rho(\R^d)$}\label{S:weightedspace}

Let $\rho$ be some nonnegative function on $\R^d$. We denote by $L^2_\rho(\R^d)$ the weighted
Hilbert space with the inner produce $\langle f,g \rangle_\rho := \int_{\R^d}f(x)g(x)\rho(x) \: \ud
x$ and set $\Norm{f}_\rho^2 :=\int_{\R^d}f(x)^2\rho(x) \: \ud x$. Following Tessitore and Zabczyk
\cite{tessitore.zabczyk:98:invariant}, we make the following definition:

\begin{definition} \label{D:Rho}
  $\rho:\R^d\mapsto \R$ is called an \textit{admissible weight} if $\rho$ is a strictly positive,
  bounded and continuous function in $L^1(\R^d)$ such that for $T>0$ there exists a constant
  $C_\rho(T)$ such that
  \begin{equation} \label{E:aw}
    \big(G(t,\cdot)*\rho(\cdot)\big)(x) \le C_\rho(T) \rho(x)
    \quad \text{for all $t\in[0,T]$ and $x\in\R^d$,}
  \end{equation}
  or equivalently,
  \begin{equation} \label{E:awbd1}
    \sup_{x \in \R^d,\: t\in[0,T]}\frac{\left(G(t,\cdot)*\rho(\cdot)\right)(x)}{\rho(x)}
    \le C_{\rho}(T).
  \end{equation}
\end{definition}

Here are several examples of admissible weights (see Example in Section 2 of
\cite{tessitore.zabczyk:98:invariant})
\begin{equation} \label{E:admRho}
  \begin{aligned}
    \rho(x) & = \exp(-a|x|), \quad               &  & a>0 \quad \text{and} \\
    \rho(x) & = \left(1+|x|^a\right)^{-1}, \quad &  & a>d.
  \end{aligned}
\end{equation}

\begin{remark} \label{R:admRho}
  The smaller the weight function $\rho(\cdot)$ (not necessarily to be admissible here) is, the
  larger the space $L_\rho^2(\R^d)$ is. For example, one may choose $\rho$ to be either a
  nonnegative function with compact support or the heat kernel itself $G(1,\cdot)$. In both cases,
  $\rho$ is smaller than those in \eqref{E:admRho} (up to a constant). However, one can easily check
  that the admissible condition \eqref{E:aw} excludes these two cases. It is much less obvious that
  the following weight function is not admissible:
  \begin{equation*}
    \rho(x) = \exp\left(-a |x|^b\right), \quad x\in\R^d,
    \quad \text{with $a>0$ and $b\in (1,2)$ fixed.}
  \end{equation*}
\end{remark}

\begin{proposition}[Proposition 2.1 of \cite{tessitore.zabczyk:98:invariant}] \label{P:weight}
  For any admissible weight $\rho$, the operators on $L^2_\rho(\R^d)$ defined by $\varphi \mapsto
  \big(G(t,\cdot)*\varphi(\cdot)\big)(x)$ can be extended to a $C_0-semigroup$ on $L^2_\rho$.
  Moreover, if $\widehat{\rho}$ is another admissible weight such that
  \begin{align*}
    \int_{\R^d}\dfrac{\rho(x)}{\widehat{\rho}(x)}\ud x < \infty,
  \end{align*}
  then for any $t<0$, the operators defined above are compact from $L^2_{\widehat{\rho}}(\R^d)$ to
  $L^2_\rho(\R^d)$.
\end{proposition}

%\textcolor{blue}{We need Lemma \ref{L:RhoDom} to how that $\varphi \in L^2_{\rho}$ will satisfy \eqref{A:icon}.}

\textcolor{magenta}{Find out if the following lemma has been used, if not we
may remove the Lemma.}

\begin{lemma}\label{L:RhoDom}
  For any admissible weight function $\rho$ and $a>0$, there exists a constant $K_{\rho,a}>0$ such
  that
  \begin{align} \label{E:RhoDom}
    \exp\left(-a|x|^2\right)\le K_{\rho,a}\: \rho(x),\qquad \text{for all $x\in\R^d$.}
  \end{align}
\end{lemma}
\begin{proof}
  Let $\rho$ be an arbitrary admissible weight function and $a>0$. By \eqref{E:aw},
  \begin{align*}
        \frac{\exp\left(-a|x|^2\right)}{\rho(x)}
    & = \frac{\exp\left(-a|x|^2\right)}{\rho(x)} \frac{\left(G(1,\cdot)*\rho(\cdot)\right)(x)}{\left(G(1,\cdot)*\rho(\cdot)\right)(x)}\\
    & \le \left(\sup_{x \in \R^d,\: t\in[0,1]}\frac{\left(G(t,\cdot)*\rho(\cdot)\right)(x)}{\rho(x)} \right)
        \frac{\exp\left(-a|x|^2\right)}{\left(G(1,\cdot)*\rho(\cdot)\right)(x)}\\
    & \le C_\rho(1) \frac{\exp\left(-a|x|^2\right)}{\left(G(1,\cdot)*\rho(\cdot)\right)(x)}.
  \end{align*}
  Noticing that for some constant $\Theta>0$, $\rho(x) \ge \Theta \prod_{\ell=1}^{d}
  \one_{[-1,1]}(x_\ell)$, a simple exercise shows that
  \begin{align*}
    C_{\rho,a}
    :=  \sup_{x\in\R^d}\frac{\exp\left(-a|x|^2\right)}{\left(G(1,\cdot)*\rho(\cdot)\right)(x)}
    \le \frac{1}{\Theta} \left(\sup_{x\in\R}\frac{\exp\left(-a x^2\right)}{\left(G(1,\cdot)*\one_{[-1,1]}(\cdot)\right)(x)}\right)^d
    <   \infty.
  \end{align*}
  This proves the lemma with $K_{\rho,a} = C_\rho(1) C_{\rho,a}$.
\end{proof}

\begin{theorem}[Theorem 1.2 \cite{chen.huang:19:comparison}] \label{T:Solve}
Suppose that
  \begin{enumerate}[(i)]
   \item the initial measure $\mu$ satisfies the following integrability condition:
   \begin{equation} \label{A:icon}
   \int_{\R^d}\exp\left(-a|x|^2\right)\Norm{\mu(\ud x)}_2 < \infty \quad \text{for all $a>0$},
   \end{equation}
   \item the spectral measure $\widehat{f}$ satisfies {\it Dalang's condition}, namely,
   \begin{equation} \label{A:dalcon}
   \Upsilon(\beta) := (2\pi)^{-d}\int_{\R^d}\dfrac{\widehat{f}(\ud \xi)}{\beta + |\xi|^2} < \infty
   \quad \text{for some and therefore any $\beta>0$}.
   \end{equation}
  \end{enumerate}
then \eqref{E:SHE} has a unique random field solution starting from $\mu$. Moreover, the solution is $L^2(\Omega)$ continuous and is adapted to the filtration induced by ${W}$.
\end{theorem}

\begin{theorem}[Theorem 1.7 \cite{chen.huang:19:comparison}] \label{T:ExUn}
  Under the assumptions of Theorem \ref{T:Solve}, for any $t>0$ and $x \in
  \R^d$, the solution to \eqref{E:SHE}, $u(t,x)$, given by \eqref{E:gmsol} is
  in $L^p(\Omega)$ and for $p\ge 2$
  \begin{equation} \label{E:mb}
    \Norm{u(t,x)}_p \le \big[\bar{\varsigma}+\sqrt{2}(G(t,\cdot) * |\mu|)(x)\big]H(t;\gamma_p)^{1/2},
  \end{equation}
  where $\bar{\varsigma}=L_0/L_b$, $\gamma_p = 32pL_b^2$ (see \eqref{E:lipcon}
  for $L_0$ and $L_b$) and the function $H(t;\gamma_p)$ is nondecreasing in $t$
  (see \cite{chen.huang:19:comparison} for the expression of $H$). Moreover, if
  \begin{equation}\label{A:dalcon0}
    \Upsilon(0) < \infty \quad \text{and} \quad \gamma_p < \frac{1}{4\Upsilon(0)},
  \end{equation}
  then this function $H(t;\gamma_p)$ is bounded in $t$. Hence, by choosing
  $C=\left(\sqrt{2}\vee \overline{\varsigma}\right)\sup_{t\ge 0}
  H(t;\gamma_p)^{1/2}<\infty$, we have that
  \begin{equation} \label{E:mb0}
    \Norm{u(t,x)}_p \le C\bigg(1+ (G(t, \cdot)*|\mu|)(x) \bigg).
  \end{equation}
\end{theorem}

\begin{corollary}
 Under the assumption \eqref{A:dalcon0} above, the solution $v = v_{t_0}$ to the restarted equation \eqref{E:restart} satisfies the following moment bound for any $q \ge 2$:
 \begin{equation}\label{E:MomBddRestart}
  \Norm{v(t,x)}_q \le C\bigg(1 + (G(t + t_0, \cdot) * |\mu|)(x)  \bigg), \quad \forall \; t \ge 0,
 \end{equation}
 where $C>0$ does not depend on $t$.
\end{corollary}
 \begin{proof}
The proof is a simple consequence of Theorem \ref{T:ExUn}. Indeed, for each $t \ge 0$
 \begin{align*}
  \Norm{v(t,x)}_q = \Norm{u(t+t_0,x)}_q \le C\bigg(1 + (G(t+t_0, \cdot)* |\mu|)(x) \bigg).
 \end{align*}
 \end{proof}

% \subsection{Invariant measure}
%
% \begin{definition}[Tightness of a family of measures] \label{D:tight}
%   Consider the Hausdorff space $(X,T)$ and suppose that the sigma algebra $\Sigma$ contains $T$. Let
%   $M$ be a collection of probability measures defined on $\Sigma$. The collection $M$ is said to be
%   \textit{tight} if for any $\epsilon>0$, there exists a compact set $K_{\epsilon} \subset X$ such
%   that for any measure $m\in M$, $m(K_{\epsilon}) > 1 - \epsilon.$
% \end{definition}

The following factorization property of the heat kernel will be used in this
paper
\begin{equation} \label{E:gsim}
  G(t,x)G(s,y)=G\left(\frac{ts}{t+s}, \frac{sx + ty}{t+s}\right)G\left(t+s, x-y\right),
\end{equation}
which can be easily verified. % See \cite[Lemma A.4]{chen.dalang:15:moments}.

\section{Proof of part (a) of Theorem \ref{T:Main}} \label{S:mainproof}
Now, for $\alpha\in (0,1)$, $t>0$ and $x\in\R^d$, define formally
\begin{align} \label{E:Ops}
  \begin{aligned}
    \left(F_{\alpha} f\right)(t,x) & := \int_0^{t} \int_{\R^d} (t-s)^{\alpha-1}G(t-s,x-y)f(s,y)\: \ud s \ud y  \quad \text{and} \\
    \left(Y_{\alpha} f\right)(t,x) & := \int_0^{t} \int_{\R^d} (t-s)^{-\alpha }G(t-s,x-y)f(s,y) W(\ud s,\ud y) \:.
  \end{aligned}
\end{align}
For $F_\alpha$, the Step 1 of the proof of \cite[Theorem
3.1]{tessitore.zabczyk:98:invariant} showed the following proposition:
\begin{proposition} \label{P:F-Comp}
  Let $\rho$ and $ \widehat{\rho}$ be given in condition (i) of Theorem
  \ref{T:Main} (see \eqref{E:rhorho}). For any $q>2$, $t_0>0$ and $\alpha \in
  (q^{-1},2^{-1})$, the operator $F_\alpha$ as an operator from
  $L^q\big((0,t_0);\: L^2_{\widehat{\rho}}\: \big)$ to $L^2_\rho(\R^d)$ is
  compact.
\end{proposition}

As for $Y_\alpha$, we have the following result:

% {{{ -- Removed stuff
%%%%%%%%%%%%%%%%%%%%%%%%%%%%%%%
%\begin{lemma} \label{L:yineq}
%  For any $\alpha\in (0,1/2)$, denote (with slightly abuse of notation)
%  \begin{equation} \label{E:Y}
%    Y_u(s,y):= \left[Y_\alpha b\left(\circ, u(\cdot,\circ)\right)\right](s,y) = \int_0^s \int_{\R^d} (s-r)^{-\alpha}G(s-r,y-z)b(z,u(r,z))W(\ud r,\ud z).
%  \end{equation}
%  Under the assumptions of Theorem \ref{T:Lrho} and (iii) of Theorem \ref{T:Main}, then for all $q\ge 2$,
%  $s>0$ and $y\in\R^d$,
%  \begin{align} \label{E:NormY}
%    \Norm{Y(s,y)}_q^2 \le C_q\: [J_q^*(s,y)]^2\: \mathcal{H}_\alpha(s).
%  \end{align}
%  Moreover, if
%  \begin{enumerate}
%    \item[(iv')] for some $q\ge 2$ and $\alpha\in \left(0,1/2\right)$, the weight function $\rho$
%      (not necessarily admissible)
%      \begin{align} \label{E:GH-int'}
% \int_{0}^t \left[\mathcal{G}_{\rho, q}(s)\:\mathcal{H}_\alpha(s)\right]^{q/2} \ud s <\infty, \quad \text{for all $t>0$},
%      \end{align}
%  \end{enumerate}
%  then there exists a constant $\Theta = \Theta\left(q,L_b,L_0,\alpha\right)$ such that
%  \begin{equation} \label{E:yineq}
%    \E \left(\int_0^{t} \Norm{Y(s,\cdot)}_{\rho}^q \ud s\right)
%    \le \Theta \int_0^t \left[\mathcal{G}_{\rho, q}(s)\: \mathcal{H}_\alpha(s)\right]^{q/2} \ud s
%    <   \infty.
%  \end{equation}
%\end{lemma}
%\begin{proof}
%  In the proof, we use $C$ to denote a generic constant that may change its value at each
%  appearance. First, an application of Minkowski's inequality shows that
%  \begin{align} \label{E:EYqr}
%    \E\left( \Norm{Y(s,\cdot)}_{\rho}^q\right)
%    = \Norm{\int_{\R^d} Y(s,y)^2\rho(y) \ud y }_{q/2}^{q/2}
%    \le \left(\int_{\R^d} \Norm{ Y(s,y) }_q^2\rho(y) \ud y \right)^{q/2}.
%    % = \Norm{\Norm{Y(s,\cdot)}_q}_\rho^q.
%  \end{align}
%  Now we look to find an appropriate upper bound for $\Norm{Y(s,y)}^2_q$. By the
%  Burkholder-Davis-Gundy inequality, as well as Minkowski's integral inequality to see that
%  \begin{align*}
%    \Norm{Y(s,y)}_q^2 \le
%      C_q \int_0^s \ud r \: (s-r)^{-2\alpha} \iint_{\R^{2d}} \ud z_1\ud z_2 \:
%            & G(s-r,y-z_1) \Norm{b\left(z_1,u(r,z_1)\right)}_q f(z_1-z_2) \\
%     \times & G(s-r,y-z_2) \Norm{b\left(z_2,u(r,z_2)\right)}_q.
%  \end{align*}
%  Note that for the Lipschitz condition in \eqref{E:lipcon}, we have that
%  \begin{align*}
%    \left|b\left(x,u\right)\right|
%    & \le \left|b\left(x,u\right) - b\left(x,0\right)\right| + \left|b\left(x,0\right)\right|
%      \le L_b |u| + L_0 \le C (1+|u|),\quad C:=L_b\vee L_0.
%  \end{align*}
%  We apply this and the moment bound \eqref{E:mb0} to $\Norm{b(z_i, u(r,z_i))}_q$ above to see that
%  \begin{align} \label{E:NormB}
%    \Norm{b(z_i, u(r,z_i))}_q
%    \le C\left(1+\Norm{u(r,z_i)}_q\right)
%    \le C\bigg(1+(G(r,\cdot) * \Norm{\mu}_q))(z_i)\bigg) = C J_q^*(r ,z_i), \quad i=1,2,
%  \end{align}
%  where $J_q^*$ is defined in \eqref{E:J*u}. Therefore,
%  \begin{align*}
%    \Norm{Y(s,y)}_q^2
%    & \le C \int_0^s \ud r \: (s-r)^{-2\alpha} \iint_{\R^{2d}} \ud z_1\ud z_2\: f(z_1-z_2) \prod_{i=1}^2  \bigg( G(s-r,y-z_i)  J_q^*(r,z_i)\bigg)     \\
%    & = C \int_0^s \ud r \: (s-r)^{-2\alpha} \iint_{\R^{2d}} \mu^*(\ud \sigma_1) \mu^*(\ud \sigma_2) \iint_{\R^{2d}} \ud z_1\ud z_2                 \\
%    & \quad \times f(z_1-z_2) \prod_{i=1}^{2} \bigg( G(s-r,y-z_i)G(r ,z_i-\sigma_i) \bigg)                                                           \\
%    & = C \int_0^s \ud r \: (s-r)^{-2\alpha} \iint_{\R^{2d}} \mu^*(\ud \sigma_1) \mu^*(\ud \sigma_2)\: G(s ,y-\sigma_1)G(s, y-\sigma_2)               \\
%    & \quad \times\iint_{\R^{2d}} \ud z_1\ud z_2 f(z_1-z_2)  \prod_{i=1}^2 G\left(\frac{r(s-r)}{s}, z_i-\sigma_i  -\frac{r}{s}(y - \sigma_i) \right)    \\
%    & \le C (2\pi)^{-2d} \int_0^s \ud r \: (s-r)^{-2\alpha} \iint_{\R^{2d}}\mu^*(\ud \sigma_1) \mu^*(\ud \sigma_2)\: G(s ,y-\sigma_1)G(s , y-\sigma_2) \\
%    & \quad \times \int_{\R^d} \widehat{f}(\ud \xi) \exp\left(-\frac{r(s-r)}{s}|\xi|^2\right),
%  \end{align*}
%  where we have applied \eqref{E:gsim} and Plancherel's theorem. Therefore,
%  \begin{align} \label{E_:NormY}
%    \Norm{Y(s,y)}_q^2 \le C (2\pi)^{-2d} [J_q^*(s + t_0,y)]^2 \int_0^s \ud r \: (s-r)^{-2\alpha} \int_{\R^d} \widehat{f}(\ud \xi)\exp\left(-\frac{r (s-r)}{s}|\xi|^2\right).
%  \end{align}
%  By symmetry of $r(s-r)/s$ term and the fact that $r(s-r)/s \ge r/2$ for all $r\in [0,s/2]$, we see
%  that the above double integral is bounded by
%  \begin{align*}
%    \le & \: 2\int_0^{s/2} \ud r \: r^{-2\alpha}\int_{\R^d} \widehat{f}(\ud \xi)\exp\left(-\frac{r}{2}|\xi|^2\right)      \\
%    =   & \: 2^{2(1-\alpha)} \int_0^{s/4} \ud r \: r^{-2\alpha}\int_{\R^d} \widehat{f}(\ud \xi)\exp\left(-r|\xi|^2\right) \\
%    \le & \: 2^{2(1-\alpha)} \int_0^{s} \ud r \: r^{-2\alpha}\int_{\R^d} \widehat{f}(\ud \xi)\exp\left(-r|\xi|^2\right)   \\
%    =   & \: 2^{2(1-\alpha)} \mathcal{H}_\alpha(s).
%  \end{align*}
%  Plugging the above bound back to \eqref{E_:NormY} proves \eqref{E:NormY}. \bigskip
%
%
%  It remains to prove \eqref{E:yineq}. By the definition of $\mathcal{G}_{\rho}(t)$ in \eqref{E:G}
%  and by \eqref{E:NormY}, we see that
%  \begin{align*}
%    \int_{\R^d}\Norm{Y(s,y)}_q^2 \: \rho(y)\ud y
%    & \le C \: \mathcal{G}_{\rho, q}(s) \mathcal{H}_\alpha(s).
%  \end{align*}
%  Plugging the above expression to the far right-hand of \eqref{E:EYqr} proves \eqref{E:yineq}.
%  Finally, the finiteness of the upper bound in \eqref{E:yineq} is guaranteed by condition
%  \eqref{E:GH-int'}. This completes the proof of Lemma \ref{L:yineq}.
%\end{proof}
%%%%%%%%%%%%%%%%%%%%%%%%%%
% }}}

\begin{lemma} \label{L:yineq}
  For any $\alpha\in (0,1/2)$, denote (with slight abuse of notation)
  \begin{equation} \label{E:Y}
    Y_{v_{t_0}}(s,t)  = Y_v(s,y)
                     := \left[Y_\alpha b\left(\circ, v(\cdot,\circ)\right)\right](s,y)
                      = \int_0^s \int_{\R^d} (s-r)^{-\alpha}G(s-r,y-z)b(z,v(r,z))W_{t_0}(\ud r,\ud z).
  \end{equation}
  Under the assumptions of Theorem \ref{T:Lrho} and (iii) of Theorem
  \ref{T:Main}, then for all $q\ge 2$, $s>0$ and $y\in\R^d$,
  \begin{align} \label{E:NormY}
  \Norm{Y_v(s,y)}_q^2 \le C_q\: [J_u^*(s + t_0,y)]^2\: \mathcal{H}_\alpha(s).
  \end{align}
  Moreover, if
  \begin{enumerate}
    \item[(iv')] for some $q\ge 2$ and $\alpha\in \left(0,1/2\right)$, the
      weight function $\rho$ (not necessarily admissible)
    \begin{align} \label{E:GH-int'}
      \int_{0}^t \left[\mathcal{G}^u_{\rho}(s + t_0)\:\mathcal{H}_\alpha(s)\right]^{q/2} \ud s <\infty, \quad \text{for all $t>0$},
    \end{align}
  \end{enumerate}
  then there exists a constant $\Theta = \Theta\left(q,L_b,L_0,\alpha\right)$
  such that
  \begin{equation} \label{E:yineq}
    \E \left(\int_0^{t} \Norm{Y_v(s,\cdot)}_{\rho}^q \ud s\right)
    \le \Theta \int_0^t \left[\mathcal{G}^u_{\rho}(s + t_0)\: \mathcal{H}_\alpha(s)\right]^{q/2} \ud s
    <   \infty.
  \end{equation}
\end{lemma}
\begin{proof}
  In the proof, we use $C$ to denote a generic constant that may change its
  value at each appearance. First, an application of Minkowski's inequality
  shows that
  \begin{align} \label{E:EYqr}
    \E\left( \Norm{Y_v(s,\cdot)}_{\rho}^q\right)
    = \Norm{\int_{\R^d} Y_v(s,y)^2\rho(y) \ud y }_{q/2}^{q/2}
    \le \left(\int_{\R^d} \Norm{ Y_v(s,y) }_q^2\rho(y) \ud y \right)^{q/2}.
    % = \Norm{\Norm{Y(s,\cdot)}_q}_\rho^q.
  \end{align}
  Now we look to find an appropriate upper bound for $\Norm{Y(s,y)}^2_q$. By the
  Burkholder-Davis-Gundy inequality, as well as Minkowski's integral inequality
  to see that
  \begin{align*}
    \Norm{Y_v(s,y)}_q^2 \le
    C_q \int_0^s \ud r \: (s-r)^{-2\alpha} \iint_{\R^{2d}} \ud z_1\ud z_2 \:
           & G(s-r,y-z_1) \Norm{b\left(z_1,v(r,z_1)\right)}_q f(z_1-z_2) \\
    \times & G(s-r,y-z_2) \Norm{b\left(z_2,v(r,z_2)\right)}_q.
  \end{align*}
  Note that for the Lipschitz condition in \eqref{E:lipcon}, we have that
  \begin{align*}
    \left|b\left(x,v\right)\right|
    & \le \left|b\left(x,v\right) - b\left(x,0\right)\right| + \left|b\left(x,0\right)\right|
    \le L_b |v| + L_0 \le C (1+|v|),\quad C:=L_b\vee L_0.
  \end{align*}
  We apply this and the moment bound \eqref{E:MomBddRestart} to $\Norm{b(z_i,
  v(r,z_i))}_q$ above to see that
  \begin{align} \label{E:NormB}
    \Norm{b(z_i, v(r,z_i))}_q
    \le C\left(1+\Norm{v(r,z_i)}_q\right)
    \le C\bigg(1+(G(r+ t_0,\cdot) * |\mu|))(z_i)\bigg) = C J_u^*(r + t_0,z_i), \quad i=1,2,
  \end{align}
  where $J_u^*$ is defined in \eqref{E:J*u}. Therefore,
  \begin{align*}
    \Norm{Y_v(s,y)}_q^2
    & \le C \int_0^s \ud r \: (s-r)^{-2\alpha} \iint_{\R^{2d}} \ud z_1\ud z_2\: f(z_1-z_2) \prod_{i=1}^2  \bigg( G(s-r,y-z_i)  J_u^*(r + t_0,z_i)\bigg)                        \\
    & = C \int_0^s \ud r \: (s-r)^{-2\alpha} \iint_{\R^{2d}} \mu^*(\ud \sigma_1) \mu^*(\ud \sigma_2) \iint_{\R^{2d}} \ud z_1\ud z_2                                            \\
    & \quad \times f(z_1-z_2) \prod_{i=1}^{2} \bigg( G(s-r,y-z_i)G(r + t_0,z_i-\sigma_i) \bigg)                                                                                \\
    & = C \int_0^s \ud r \: (s-r)^{-2\alpha} \iint_{\R^{2d}} \mu^*(\ud \sigma_1) \mu^*(\ud \sigma_2)\: G(s + t_0,y-\sigma_1)G(s + t_0,y-\sigma_2)                              \\
    & \quad \times\iint_{\R^{2d}} \ud z_1\ud z_2 f(z_1-z_2)  \prod_{i=1}^2 G\left(\frac{(r+t_0)(s-r)}{s + t_0}, z_i-\sigma_i\frac{r + t_0}{s + t_0}-\frac{s-r}{s+t_0}y \right) \\
    & \le C (2\pi)^{-2d} \int_0^s \ud r \: (s-r)^{-2\alpha} \iint_{\R^{2d}}\mu^*(\ud \sigma_1) \mu^*(\ud \sigma_2)\: G(s + t_0,y-\sigma_1)G(s + t_0, y-\sigma_2)               \\
    & \quad \times \int_{\R^d} \widehat{f}(\ud \xi) \exp\left(-\frac{(r + t_0)(s-r)}{s+ t_0}|\xi|^2\right),
  \end{align*}
  where we have applied \eqref{E:gsim} and Plancherel's theorem. Therefore,
  \begin{align*}
   \Norm{Y_v(s,y)}_q^2 \le C (2\pi)^{-2d} [J_u^*(s + t_0,y)]^2 \int_0^s \ud r \: (s-r)^{-2\alpha} \int_{\R^d} \widehat{f}(\ud \xi)\exp\left(-\frac{(r + t_0)(s-r)}{s + t_0}|\xi|^2\right).
  \end{align*}
  Notice the function
  \[
   t_0 \mapsto \frac{r + t_0}{s + t_0} \quad \text{for } t_0 > 0,
  \]
  is increasing in $t_0 $ whenever $s > r > 0$. Thus we have that
  \begin{equation} \label{E_:NormY}
    \Norm{Y_v(s,y)}_q^2 \le C (2\pi)^{-2d} [J_u^*(s + t_0,y)]^2 \int_0^s \ud r \: (s-r)^{-2\alpha} \int_{\R^d} \widehat{f}(\ud \xi)\exp\left(-\frac{r(s-r)}{s}|\xi|^2\right).
  \end{equation}
  Furthermore, by symmetry of $r(s-r)/s$ and the fact that $r(s-r)/s \ge r/2$
  for all $r\in [0,s/2]$, we see that the above double integral is bounded by
  \begin{align*}
    \le & \: 2\int_0^{s/2} \ud r \: r^{-2\alpha}\int_{\R^d} \widehat{f}(\ud \xi)\exp\left(-\frac{r}{2}|\xi|^2\right)      \\
    =   & \: 2^{2(1-\alpha)} \int_0^{s/4} \ud r \: r^{-2\alpha}\int_{\R^d} \widehat{f}(\ud \xi)\exp\left(-r|\xi|^2\right) \\
    \le & \: 2^{2(1-\alpha)} \int_0^{s} \ud r \: r^{-2\alpha}\int_{\R^d} \widehat{f}(\ud \xi)\exp\left(-r|\xi|^2\right)   \\
    =   & \: 2^{2(1-\alpha)} \mathcal{H}_\alpha(s).
  \end{align*}
  Plugging the above bound back to \eqref{E_:NormY} proves \eqref{E:NormY}.
  \bigskip

  It remains to prove \eqref{E:yineq}. By the definition of $\mathcal{G}^u_{\rho}(t)$ in \eqref{E:G}
  and by \eqref{E:NormY}, we see that
  \begin{align*}
    \int_{\R^d}\Norm{Y_v(s,y)}_q^2 \: \rho(y)\ud y
    & \le C \: \mathcal{G}^u_{\rho}(s + t_0) \mathcal{H}_\alpha(s).
  \end{align*}
  Plugging the above expression to the far right-hand of \eqref{E:EYqr} proves \eqref{E:yineq}.
  Finally, the finiteness of the upper bound in \eqref{E:yineq} is guaranteed by condition
  \eqref{E:GH-int'}. This completes the proof of Lemma \ref{L:yineq}.
\end{proof}

\begin{lemma}[Factorization lemma] \label{L:Factorize}
  Let $u(t ,x ,\mu)$ be the solution to \eqref{E:SHE} and $Y_v$ be as in
  \eqref{E:Y}. Under the conditions of Theorem \ref{T:Lrho} and (iii) of
  Theorem \ref{T:Main}, for all $\alpha \in (0,1/2)$, the following
  factorization holds
  \begin{equation*}
    \dfrac{\sin(\alpha \pi)}{\pi}
    \int_0^t(t-s)^{\alpha-1}\left[G(t-s,\cdot)*Y_v(s,\cdot)\right](x) \ud s
    = \int_0^t  \int_{\R^d} G(t-r,x-z)b(z,v(r,z)) W_{t_0}(\ud r,\ud z).
  \end{equation*}
  As a consequence,
  \begin{equation} \label{E:1fac}
    v(t,x)= [G(t, \cdot) * u(t_0,\cdot, \mu)](x) +\dfrac{\sin(\alpha \pi)}{\pi} \left[F_\alpha Y_v\right](t,x),
    \quad \text{for all $t>0$ and  $x\in\R^d$}.
  \end{equation}
\end{lemma}
\begin{proof}
  The lemma is straightforward provided that one can switch the orders of
  stochastic and ordinary integrals:
  \begin{align}
      & \int_0^t(t-s)^{\alpha-1}\left[G(t-s,\cdot)*Y_v(s,\cdot)\right](x) \ud s \notag                                                        \\
    = & \int_0^t \ud s (t-s)^{\alpha-1}\int_{\R^d} \ud y\: G(t-s,x-y) \notag                                                                  \\
      & \times \int_0^s\: \int_{\R^d} \: (s-r)^{-\alpha}G(s-r,y-z)b(z,v(r,z))W_{t_0}(\ud r,\ud z) \notag                                      \\
    = & \int_0^t \ud s\: (t-s)^{\alpha-1} \int_0^s \int_{\R^d} \: (s-r)^{-\alpha}G(t-r,x-z)b(z,v(r,z))W_{t_0}(\ud r,\ud z) \label{E:Fubini_1} \\
    = & \int_0^t  \int_{\R^d} W(\ud r,\ud z)G(t-r,x-z)b(v(r,z)) \int_r^t \ud s\: (s-r)^{-\alpha} (t-s)^{\alpha-1} \label{E:Fubini_2}          \\
    = & \dfrac{\pi}{\sin(\alpha \pi)}\int_0^t  \int_{\R^d} G(t-r,x-z)b(v(r,z)) W_{t_0}(\ud r,\ud z),\notag
  \end{align}
  where the last step is the Beta integral which requires that $\alpha\in
  (0,1)$. It remains to justify the two applications of the stochastic Fubini's
  theorem (see Theorem 5.30 of Chapter one in
  \cite{dalang.khoshnevisan.ea:09:minicourse}, or also
  \cite{walsh:86:introduction} or Theorem 4.18 of
  \cite{da-prato.zabczyk:92:stochastic}) in \eqref{E:Fubini_1} and
  \eqref{E:Fubini_2} in the following two steps. \bigskip

  {\bf\noindent Step 1.~} In this step, we justify the change of orders in
  \eqref{E:Fubini_1}. Note that $t, x$ and $s$ are fixed. It suffices to prove
  the following condition:
  \begin{align*}
  I_1:= & \int_{\R^d}\ud y \: G(t-s,x-y) \int_0^s \ud r\: (s-r)^{-2\alpha} \iint_{\R^{2d}} \ud z_1\ud z_2\:               \\
        & \times f(z_1-z_2) \left(  \prod_{i=1}^{2} G(s-r,y-z_i)\right)  \: \E\left(\prod_{i=1}^{2}b(z_i,v(r,z_i))\right) \\
     =  & \int_{\R^d}\ud y \: G(t-s,x-y) \Norm{Y_v(s,y)}_2^2                                                              \\
     <  & \: +\infty.
  \end{align*}
  But this follows immediately from \eqref{E:NormY}. Indeed
  \begin{align*}
    \int_{\R^d}\ud y \: G(t-s,x-y) \Norm{Y_v(s,y)}_2^2
    & \le C \int_{\R^d} \ud y \;  G(t-s,x-y)\left(J_u^*(s+t_0,y)\right)^2 \mathcal{H}_\alpha(s) \\
    & =C \mathcal{H}_\alpha(s) \int_{\R^d} \ud y \;  G(t-s,x-y)\left(J_u^*(s + t_0,y)\right)^2.
  \end{align*}
  Then Condition \eqref{E:G} implies that $I_1 < \infty$ as $\rho(y)$ cannot
  decay as fast as $y \mapsto G(t-s,x-y)$ (See Remark \ref{R:admRho}).

  {\bf\noindent Step 2.~} Similarly, as for \eqref{E:Fubini_2}, we need to show
  that
  \begin{align*}
    I_2:= & \int_0^t \ud s \: (t-s)^{\alpha-1} \int_0^s \ud r\: (s-r)^{-2\alpha} \iint_{\R^{2d}} \ud z_1\ud z_2\: \\
          & \times f(z_1-z_2) \left(  \prod_{i=1}^{2} G(t-r,x-z_i)\right)  \: \E\left(\prod_{i=1}^{2}b(z_i,v(r,z_i))\right) < \infty.
  \end{align*}
  By the Cauchy Schwartz inequality, \eqref{E:NormB} and because $\alpha\in
  (0,1/2)$,
  \begin{align*}
    I_2\le & C \int_0^t \ud s \: (t-s)^{\alpha-1} \int_0^s \ud r\: (s-r)^{-2\alpha} \iint_{\R^{2d}} \ud z_1\ud z_2\:f(z_1-z_2) \prod_{i=1}^{2}\bigg(  G(t-r,x-z_i) J_u^*(r + t_0,z_i)\bigg) \\
         = & C' \int_0^t \ud r \: (t-r)^{-\alpha}  \iint_{\R^{2d}} \ud z_1\ud z_2\:f(z_1-z_2) \prod_{i=1}^{2}\bigg(  G(t-r,x-z_i) J_u^*(r + t_0, z_i)\bigg).
  \end{align*}
  Now by the same arguments as those leading to \eqref{E:NormY} (with $2\alpha$
  there replaced by $\alpha$), we see that
  \begin{align*}
    I_2 & \le C (J_u^*(t,x))^2 \int_0^t \ud r \: r^{-\alpha} \int_{\R^d} \widehat{f}(\ud \xi)\exp\left(-\frac{r(t-r)}{t}|\xi|^2\right),
  \end{align*}
  which is finite by \eqref{E:H} where we replace $\alpha$ with $\alpha/2$ and
  repeat the same steps right after Equation \eqref{E_:NormY}. This completes
  the proof of Lemma \ref{L:Factorize}.
\end{proof}

% {{{ -- Removed stuff
%\begin{lemma}[Factorization lemma] \label{L:Factorize}
%  Let $u(t,x)$ be the solution to \eqref{E:SHE} and $Y$ be given in \eqref{E:Y}. Under the
%  conditions of Theorem \ref{T:Lrho} and (iii) of Theorem \ref{T:Main}, for all $\alpha \in
%  (0,1/2)$, the following factorization holds
%  \begin{equation*}
%      \dfrac{\sin(\alpha \pi)}{\pi}
%      \int_0^t(t-s)^{\alpha-1}\left[G(t-s,\cdot)*Y(s,\cdot)\right](x) \ud s
%    = \int_0^t  \int_{\R^d} G(t-r,x-z)b(z,u(r,z)) W(\ud r,\ud z).
%  \end{equation*}
%  As a consequence,
%  \begin{equation} \label{E:1fac}
%    u(t,x)=J_0(t,x)+\dfrac{\sin(\alpha \pi)}{\pi} \left[F_\alpha Y\right](t,x),
%    \quad \text{for all $t>0$ and  $x\in\R^d$}.
%  \end{equation}
%\end{lemma}
%\begin{proof}
%  % \begin{proof}[Sketch of proof.]
%  The lemma is straightforward provided that one can switch the orders of stochastic and ordinary
%  integrals:
%  \begin{align}
%      & \int_0^t(t-s)^{\alpha-1}\left[G(t-s,\cdot)*Y(s,\cdot)\right](x) \ud s                                        \notag             \\
%    = & \int_0^t \ud s (t-s)^{\alpha-1}\int_{\R^d} \ud y\: G(t-s,x-y)                                                \notag             \\
%      & \times \int_0^s\: \int_{\R^d} \: (s-r)^{-\alpha}G(s-r,y-z)b(z,u(r,z))W(\ud r,\ud z)                          \notag             \\
%    = & \int_0^t \ud s\: (t-s)^{\alpha-1} \int_0^s \int_{\R^d} \: (s-r)^{-\alpha}G(t-r,x-z)b(z,u(r,z))W(\ud r,\ud z) \label{E:Fubini_1} \\
%    = & \int_0^t  \int_{\R^d} W(\ud r,\ud z)G(t-r,x-z)b(u(r,z)) \int_r^t \ud s\: (s-r)^{-\alpha} (t-s)^{\alpha-1}    \label{E:Fubini_2} \\
%    = & \dfrac{\pi}{\sin(\alpha \pi)}\int_0^t  \int_{\R^d} G(t-r,x-z)b(u(r,z)) W(\ud r,\ud z),\notag
%  \end{align}
%  where the last step is the Beta integral which requires that $\alpha\in (0,1)$. It remains to
%  justify the two applications of the stochastic Fubini's theorem (see Theorem 5.30 of Chapter one
%  in \cite{dalang.khoshnevisan.ea:09:minicourse}, or also \cite{walsh:86:introduction} or Theorem
%  4.18 of \cite{da-prato.zabczyk:92:stochastic}) in \eqref{E:Fubini_1} and \eqref{E:Fubini_2} in the
%  following two steps. \bigskip
%
%  {\bf\noindent Step 1.~} In this step, we justify the change of orders in \eqref{E:Fubini_1}. Note that $t, x$ and $s$ are fixed.
%  It suffices to prove the following condition:
%  \begin{align*}
%    I_1:= & \int_{\R^d}\ud y \: G(t-s,x-y) \int_0^s \ud r\: (s-r)^{-2\alpha} \iint_{\R^{2d}} \ud z_1\ud z_2\:               \\
%          & \times f(z_1-z_2) \left(  \prod_{i=1}^{2} G(s-r,y-z_i)\right)  \: \E\left(\prod_{i=1}^{2}b(z_i,u(r,z_i))\right) \\
%        = & \int_{\R^d}\ud y \: G(t-s,x-y) \Norm{Y(s,y)}_2^2                                                                \\
%      <   & \: +\infty.
%  \end{align*}
%But this follows immediately from \eqref{E:NormY}. Indeed
% \begin{align*}
%  \int_{\R^d}\ud y \: G(t-s,x-y) \Norm{Y(s,y)}_2^2  &\le \int_{\R^d} \ud y \;  G(t-s,x-y)\left(J_2^*(s,y)\right)^2 \mathcal{H}_\alpha(s) \\
%  & =C \mathcal{H}_\alpha(s) \int_{\R^d} \ud y \;  G(t-s,x-y)\left(J_2^*(s,y)\right)^2.
% \end{align*}
% Then Condition \eqref{E:G} implies that $I_1 < \infty$ as $\rho(y)$ cannot decay as fast as $y \mapsto G(t-s,x-y)$ (See Remark \ref{R:admRho}).
%
%% We can simplify the following by applying \eqref{E:NormY}.
%% See above few lines.
%%
%%  By the moment estimate in \eqref{E:NormY} with $q=2$, we see that
%%  \begin{align*}
%%    I_1 \le & C\int_{\R^d}\ud y \: G(t-s,x-y) \int_0^s \ud r \: (s-r)^{-2\alpha} \iint_{\R^{2d}} \mu_*(\ud \sigma_1) \mu_*(\ud \sigma_2) \\
%%            & \times\int_{\R^d} \widehat{f}(\ud\xi) \exp\left(-\frac{r(s-r)}{s}|\xi|^2\right)
%%   \prod_{i=1}^{2}G(s,y-\sigma_i).
%%  \end{align*}
%%  Now we bound the three heat kernels using \eqref{E:gsim} as follows:
%%  \begin{align*}
%%     G & (t-s,x-y)  \prod_{i=1}^{2}G(s,y-\sigma_i)                                                                   \\
%%       & =\frac{G\left(2(t-s),x-y\right)^2}{G\left(4(t-s,0\right)} \prod_{i=1}^{2}G\left(s,y-\sigma_i\right)         \\
%%       & \le 2^d \frac{G\left(2(t-s),x-y\right)^2}{G\left(4(t-s,0\right)} \prod_{i=1}^{2}G\left(2s,y-\sigma_i\right) \\
%%       & = 2^d [4(t-s)]^{d/2} \prod_{i=1}^{2}\bigg[ G\left(2s,y-\sigma_i\right)G\left(2(t-s),x-y\right)\bigg]        \\
%%       & = 2^{2d} (t-s)^{d/2} \prod_{i=1}^{2}\left[G\left(2t,x-\sigma_i\right)G\left(\frac{2(t-s)s}{t},y-\frac{s}{t}(x-\sigma_i)\right)\right].
%%  \end{align*}
%%  By carrying out the integration in the following order, we see that
%%  \begin{align*}
%%        & \iint_{\R^{2d}} \mu_*(\ud\sigma_1) \mu_*(\ud\sigma_2) \int_{\R^d}\ud y\: G(t-s,x-y) \prod_{i=1}^{2}G(s,y-\sigma_i)                                                                               \\
%%    \le & 2^{2d} (t-s)^{d/2} \iint_{\R^{2d}} \mu_*(\ud\sigma_1) \mu_*(\ud\sigma_2) \left( \prod_{i=1}^{2}G\left(2t,x-\sigma_i\right)\right) G\left(\frac{4(t-s)s}{t},\frac{s}{t}(\sigma_1-\sigma_2)\right) \\
%%    \le & 2^{2d} (t-s)^{d/2} \iint_{\R^{2d}} \mu_*(\ud\sigma_1) \mu_*(\ud\sigma_2) \left( \prod_{i=1}^{2}G\left(2t,x-\sigma_i\right)\right) G\left(\frac{4(t-s)s}{t},0\right)                              \\
%%    =   & C_{t,s}\: J_*^2(2t,x).
%%  \end{align*}
%%  Therefore,
%%  \begin{align*}
%%    I_1 & \le C_{t,s}\: J_*^2(2t,x) \int_0^s \ud r\: (s-r)^{-2\alpha}\int_{\R^d} \widehat{f}(\ud\xi) \exp\left(-\frac{r(s-r)}{s}|\xi|^2\right)
%%   \le C_{t,s}'\: J_*^2(2t,x) \mathcal{H}(s),
%%  \end{align*}
%%  where we have applied the same arguments as those in the proof of Lemma \ref{L:yineq} in the
%%  second inequality. Condition \eqref{E:GH-int} implies that $I_1<\infty$. \bigskip
%%\textcolor{blue}{
%%A type has been made above when using \eqref{E:NormY} to bound $\Norm{Y(s,y)}_2^2$. Note that we need to show
%% \begin{align*}
%%  \int_{\R^d} \ud y \; G(t-s,x-y) \Norm{Y(s,y)}_2^2 &\le C \int_{\R^d} \ud y \;  G(t-s,x-y)\left(J_2^*(s,y)\right)^2 \mathcal{H}_\alpha(s) \\
%%  &= \mathcal{H}_\alpha(s) \int_{\R^d} \ud y \;  G(t-s,x-y)\left(J_2^*(s,y)\right)^2 < \infty
%% \end{align*}
%%However, because $s$ and $t$ are fixed, this will follow from the assumption \eqref{E:G} which says that $\mathcal{G}_{\rho, 2}(s) < \infty$ for each $s >0$. This is due to the fact that $\exp(-|x|^2) \le C \rho(x)$.
%%}
%  {\bf\noindent Step 2.~} Similarly, as for \eqref{E:Fubini_2}, we need to show that
%  \begin{align*}
%    I_2:= & \int_0^t \ud s \: (t-s)^{\alpha-1} \int_0^s \ud r\: (s-r)^{-2\alpha} \iint_{\R^{2d}} \ud z_1\ud z_2\: \\
%          & \times f(z_1-z_2) \left(  \prod_{i=1}^{2} G(t-r,x-z_i)\right)  \: \E\left(\prod_{i=1}^{2}b(z_i,u(r,z_i))\right) < \infty.
%  \end{align*}
%  By the Cauchy Schwartz inequality, \eqref{E:NormB} and because $\alpha\in (0,1/2)$,
%  \begin{align*}
%    I_2\le & C \int_0^t \ud s \: (t-s)^{\alpha-1} \int_0^s \ud r\: (s-r)^{-2\alpha} \iint_{\R^{2d}} \ud z_1\ud z_2\:f(z_1-z_2) \prod_{i=1}^{2}\bigg(  G(t-r,x-z_i) J_2^*(r,z_i)\bigg) \\
%        =  & C' \int_0^t \ud r \: (t-r)^{-\alpha}  \iint_{\R^{2d}} \ud z_1\ud z_2\:f(z_1-z_2) \prod_{i=1}^{2}\bigg(  G(t-r,x-z_i) J_2^*(r,z_i)\bigg).
%  \end{align*}
%  Now by the same arguments as those leading to \eqref{E:NormY} (with $2\alpha$ there replaced by
%  $\alpha$), we see that
%  \begin{align*}
%    I_2 & \le C J_*^2(t,x) \int_0^t \ud r \: r^{-\alpha} \int_{\R^d} \widehat{f}(\ud \xi)\exp\left(-\frac{r(t-r)}{r}|\xi|^2\right),
%  \end{align*}
%  which is finite by \eqref{E:H} where we replace $\alpha$ with $\alpha/2$. Indeed, since $r(t-r) \ge tr/2$ for $r \in [0,t]$, then
% \begin{align*}
%  \int_0^t \ud r \: r^{-\alpha} \int_{\R^d} \widehat{f}(\ud \xi)\exp\left(-\frac{r(t-r)}{r}|\xi|^2\right) &\le \int_0^t \ud r \: r^{-\alpha} \int_{\R^d} \widehat{f}(\ud \xi)\exp\left(-\frac{r}{2}|\xi|^2\right) < \infty.
% \end{align*}
%This completes the proof of Lemma \ref{L:Factorize}.
%\end{proof}
% }}}
\begin{codes}% Mathematica

Integrate[(s-r)^{-a} (t-s)^{a-1},{s,r,t},Assumptions->{t>s>0}]

Integrate[(t-s)^{a-1} (s-r)^{-2a},{s,r,t},Assumptions->{t>r>0,0<a<1/2}]

\end{codes}

% {{{ -- Removed stuff
% I do not think we need the following lemma anymore.
%
%\begin{lemma} \label{L:Restart}
%  Let $u(t,x)$ be the solution to \eqref{E:SHE}. Under conditions (i) -- (iii) of Theorem
%  \ref{T:Lrho}, for any $q\ge 2$, there exists some constant $C$ depending on $q$ and the value of
%  $\Upsilon(0)$ such that
%  \begin{align*}
%        \E\left(\Norm{u(t,\cdot)*G(t_0,\cdot)}_{\rho}^q\right)
%    \le C\: \mathcal{G}_{\rho,\mu_*}^{q/2} (t+t_0),
%        \quad \text{for all $t$ and $t_0>0$,}
%  \end{align*}
%  where $\mu_*$ and $ \mathcal{G}_{\rho,\mu_*}$ are defined in \eqref{E:J*} and \eqref{E:G},
%  respectively.
%\end{lemma}
%\begin{proof}
%  By the same arguments as in \eqref{E:EYqr} with $Y$ replaced by $u(t,\cdot)*G(t_0,\cdot)$ to see
%  that
%  \begin{align*}
%    \E\left(\Norm{u(t,\cdot)*G(t_0,\cdot)}_{\rho}^q\right)
%    \le \left(\int_{\R^d} \Norm{ \left(u(t,\cdot)*G(t_0,\cdot)\right)(x)}_q^2\: \rho(x) \ud x \right)^{q/2}.
%  \end{align*}
%  Now we estimate $\Norm{ \left(u(t,\cdot)*G(t_0,\cdot)\right)(x)}_q^2$. Notice that
%  \begin{align*}
%    \left(u(t,\cdot)*G(t_0,\cdot)\right)(x)
%    & = J_0(t+t_0,x) + \int_0^t\int_{\R^d} G(t+t_0-s,x-y) b(y,u(s,y))W(\ud s, \ud y),
%  \end{align*}
%  where we have applied Fubini's theorem for the $J_0\left(\cdot,\cdot\right)$ term and the
%  stochastic Fubini's theorem for the stochastic integral part, the verification of which can be
%  done in the same way as those in the step 1 of the proof of Lemma \ref{L:Factorize} with
%  $\alpha=0$. Denote the above stochastic integral by $I(t,t_0,x)$. One can adapt the arguments in
%  the proof of \eqref{E:NormY} including setting $\alpha=0$ and increasing the time integral from
%  $(0,t)$ to $(0,t+t_0)$ to see that
%  \begin{align*}
%    \Norm{I(t,t_0,x)}_q^2
%    & \le C  J_*^2\left(t+t_0,x\right)\int_0^{t+t_0} \ud s \int_{\R^d} \widehat{f}(\ud \xi) \exp\left(-\frac{s(t+t_0-s)}{t+t_0}|\xi|^2\right) \\
%    & \le 2 C J_*^2\left(t+t_0,x\right)\int_0^{t+t_0} \ud s \int_{\R^d} \widehat{f}(\ud \xi) \exp\left(-s|\xi|^2\right)                       \\
%    & \le 2 C J_*^2\left(t+t_0,x\right)\int_{\R^d} \widehat{f}(\ud \xi)\int_0^{\infty} \ud s\:  \exp\left(-s|\xi|^2\right)                    \\
%    & \le 2 C J_*^2\left(t+t_0,x\right)\int_{\R^d} |\xi|^{-2} \widehat{f}(\ud \xi)                                                            \\
%    & =   2 C J_*^2\left(t+t_0,x\right) \Upsilon(0),
%  \end{align*}
%  where we have used condition \eqref{E:U0finite}. Therefore,
%  \begin{align*}
%    \Norm{\left(u(t,\cdot)*G(t_0,\cdot)\right)(x)}_q^2 \le  C J_*^2(t+t_0,x),
%  \end{align*}
%  where we have used the fact that $J_0^2(t,x)\le J_*^2(t,x)$. Multiplying both sides of the above
%  equation by $\rho(x)$ and then integrating over $x$ prove the lemma.
%\end{proof}
% }}}

\begin{codes}% Mathematica

  Integrate[Exp[-s x],{s,0,Infinity}, Assumptions->x>0]

\end{codes}

We are now ready to prove part (a) of Theorem \ref{T:Main}.\

\begin{proof}[Proof of Theorem \ref{T:Main} (a)]
Recall that Lemma \ref{L:Factorize} implies that for any $t_0 > 0$,
 \[
  v(t_0,x) = \Big(G(t_0,\cdot) * u(t_0,\cdot, \mu)\Big)(x) +\dfrac{\sin(\alpha \pi)}{\pi} \left[F_\alpha Y_v\right](t_0,x).
 \]
Now define the set $\mathscr{K} (\Lambda)$ as
 \[
  \mathscr{K} (\Lambda) = \{ \big(G(t_0,\cdot)*y(\cdot)\big)(x)+(F_\alpha h)(x),
  \quad \text{where $\Norm{y}_{\hat\rho} \le \Lambda$  and   $\Norm{h}_{L^q(0,1,L^2_{\hat \rho})} \le \Lambda $} \}.
 \]
Then by Proposition \ref{P:F-Comp}, $\mathscr{K} (\Lambda)$ is relatively compact in $L^2_{ \rho}$ for $\Lambda >0$. Moreover, by Chebyshev's inequality, \eqref{E:uInLrho} and \eqref{E:yineq} we have
 \begin{align}
  \mathbb{P}\{ v(t_0,\cdot) \not \in \mathscr{K}(\Lambda) \}
  & \le \mathbb{P}\left[ \left(\int_{0}^{t_0} \Norm{Y_v(s,\cdot)}_{\widehat{\rho}}^q \: \ud s\right)^{1/q} > \dfrac{\pi \Lambda}{\sin(\alpha \pi )} \right] + \mathbb{P}\left[ \Norm{u(t_0,\cdot, \mu)}_{\widehat{\rho}} > \Lambda \right] \nonumber \\
  & \le \dfrac{\sin^q(\alpha \pi)}{\pi^q \Lambda^q}\E \int_{0}^{t_0} \Norm{Y_v(s,\cdot)}_{\hat \rho}^q \: \ud s + \frac{1}{\Lambda^2} \mathbb{E}\left(\Norm{u(t_0,\cdot, \mu)}_{\widehat{\rho}}^2\right) \nonumber \\
  & \le \dfrac{\sin^q(\alpha \pi)}{\pi^q \Lambda^q} \Theta\int_0^{t_0} \left(\mathcal{G}^u_{\hat\rho}(s + t_0) \mathcal{H}_{\alpha}(s)\right)^{q/2} \ud s + \frac{1}{\Lambda^2} \mathbb{E}\left(\Norm{u(t_0,\cdot, \mu}_{\widehat{\rho}}^2\right) < \infty.
  \label{E:Restartbdd}
 \end{align}
With all this said, \eqref{E:Restartbdd} implies that
 \begin{align*}
  \mathbb{P}\{ v(t_0,\cdot) \in \mathscr{K}(\Lambda) \}
  & \ge 1- \dfrac{\sin^q(\alpha \pi)}{\pi^q \Lambda^q} \left(\Theta\int_0^{t_0} \left(\mathcal{G}^u_{\hat\rho}(s + t_0) \mathcal{H}_{\alpha}(s)\right)^{q/2} \ud s\right) - \frac{1}{\Lambda^2} \mathbb{E}\left(\Norm{u(t_0,\cdot, \mu)}_{\rho}^2\right),
 \end{align*}
and we can choose $\Lambda > 0$ big enough such that
 \[
  \dfrac{\sin^q(\alpha \pi)}{\pi^q \Lambda^q} \left(\Theta\int_0^{t_0} \left(\mathcal{G}^u_{\hat\rho}(s + t_0) \mathcal{H}_{\alpha}(s)\right)^{q/2} \ud s\right) + \frac{1}{\Lambda^2} \mathbb{E}\left(\Norm{u(t_0,\cdot, \mu)}_{\rho}^2\right) < \epsilon,
 \]
thus we get that
 \begin{align*}
  \mathbb{P}\{ v(t_0,\cdot) \in \mathscr{K}(\Lambda) \} & \ge 1- \epsilon.
 \end{align*}
We can now deduce from this and the semigroup property of the solution to the heat equation that for any $t > t_0$,
 \begin{equation}\label{E:vtight}
  \mathbb{P}\{ v(t,\cdot) \in \mathscr{K}(\Lambda) \} \ge (1-
  \epsilon)\mathbb{P}\{ \Norm{v(t-t_0,\cdot)}_{\hat\rho} < \Lambda
  \}.
 \end{equation}
We now apply Chebyshev's inequality and \eqref{E:InitData} to see that for all $t > t_0$ that there exists a $C>0$ such that
 \[
  \mathbb{P}\{ \Norm{v(t-t_0,\cdot)}_{\hat\rho} \ge \Lambda
  \} \le \frac{\mathbb{E}\left( \Norm{v(t-t_0,\cdot)}_{\hat\rho}^2 \right)}{\Lambda^2} = \frac{\mathbb{E}\left( \Norm{u(t,\cdot, \mu)}_{\hat\rho}^2 \right)}{\Lambda^2} \le C\frac{\mathcal{G}^u_{\hat{\rho}}\left( t \right) }{\Lambda^2} \le \frac{C}{\Lambda^2}.
 \]
Thus, by enlarging $\Lambda$ if necessary,
\[
\mathbb{P}\{ \Norm{v(t-t_0,\cdot)}_{\hat\rho} \ge \Lambda\} = \mathbb{P}\{ \Norm{u(t,\cdot, \mu)}_{\hat\rho} \ge \Lambda
\}  \le \epsilon, \quad \forall \; t \ge t_0.
\]
Combining this with \eqref{E:vtight} yields
 \[
  \mathbb{P}\left\{ u(t, \cdot, \mu) \in \mathscr{K}(\Lambda) \right\} \ge (1-\epsilon), \quad t \ge t'_0,
 \]
where $t'_0 = 2 t_0$. Since $t_0 >0$ is arbitrary, then this completes the proof of \eqref{E:cpt}.
\end{proof}

% {{{ -- Removed Stuff
%%%%%%%%%%%%%%%%%%%%%%%%%%%
%\begin{proof}[Proof of Theorem \ref{T:Main} (b)]
% Recall that from Lemma \ref{L:Factorize} implies that for any $t_0 > 0$,
% \[
% u(t_0,x) = \Big(G(t_0,\cdot) * u(0,\cdot)\Big)(x) +\dfrac{\sin(\alpha \pi)}{\pi} \left[F_\alpha Y\right](t_0,x),
% \]
% and again recall that we assume $u(0,\cdot)$ satisfies \eqref{E:RandInit}. Now define the set $\mathscr{K} (\Lambda)$ as
% \[
% \mathscr{K} (\Lambda) = \{ \big(G(t_0,\cdot)*y(\cdot)\big)(x)+(F_\alpha h)(x),
% \quad \text{where $\Norm{y}_{\hat\rho} \le \Lambda$  and  $\Norm{h}_{L^q(0,1,L^2_{\hat \rho})} \le \Lambda $} \}.
% \]
% Then by Proposition \ref{P:F-Comp}, $\mathscr{K} (\Lambda)$ is relatively compact in $L^2_{ \rho}$ for $\Lambda >0$. Moreover, by \eqref{E:yineq} and Chebyshev's inequality we have
% \begin{align}
% \mathbb{P}\{ u(t_0,\cdot) \not \in \mathscr{K}(\Lambda) \}
% & \le \mathbb{P}\left[ \left(\int_{0}^{t_0} \Norm{Y(s,\cdot)}_{\widehat{\rho}}^q \: \ud s\right)^{1/q} > \dfrac{\pi \Lambda}{\sin(\alpha \pi )} \right] + \mathbb{P}\left[ \Norm{u(0,\cdot)}_{\widehat{\rho}} > \Lambda \right] \nonumber \\
% & \le \dfrac{\sin^q(\alpha \pi)}{\pi^q \Lambda^q}\E \int_{0}^{t_0} \Norm{Y(s,\cdot)}_{\hat \rho}^q \: \ud s + \frac{1}{\Lambda^2} \mathbb{E}\left(\Norm{u(0,\cdot}_{\widehat{\rho}}^2\right) \nonumber \\
% & \le \dfrac{\sin^q(\alpha \pi)}{\pi^q \Lambda^q} \Theta\int_0^{t_0} \left(\mathcal{G}_{\hat\rho, q}(s) \mathcal{H}_{\alpha}(s)\right)^{q/2} \ud s + \frac{1}{\Lambda^2} \mathbb{E}\left(\Norm{u(0,\cdot}_{\widehat{\rho}}^2\right) < \infty \label{E:Restartbdd}
% \end{align}
% With all this said, \eqref{E:Restartbdd} implies that
% \begin{align*}
% \mathbb{P}\{ u(t_0,\cdot) \in \mathscr{K}(\Lambda) \}
% & \ge 1- \dfrac{\sin^q(\alpha \pi)}{\pi^q \Lambda^q} \left(\Theta\int_0^{t_0} \left(\mathcal{G}_{\hat\rho, q}(s) \mathcal{H}_{\alpha}(s)\right)^{q/2} \ud s\right) - \frac{1}{\Lambda^2} \mathbb{E}\left(\Norm{u(0,\cdot)}_{\rho}^2\right)
% \end{align*}
% Therefore, we can choose $\Lambda > 0$ big enough such that
% \[
% \dfrac{\sin^q(\alpha \pi)}{\pi^q \Lambda^q} \left(\Theta\int_0^{t_0} \left(\mathcal{G}_{\hat\rho, q}(s) \mathcal{H}_{\alpha}(s)\right)^{q/2} \ud s\right) + \frac{1}{\Lambda^2} \mathbb{E}\left(\Norm{u(0,\cdot)}_{\rho}^2\right) < \epsilon,
% \]
% thus we get that
% \begin{align*}
% \mathbb{P}\{ u(t_0,\cdot) \in \mathscr{K}(\Lambda) \} & \ge 1- \epsilon.
% \end{align*}
% We can now deduce from this and the semigroup property of the solution to the heat equation that for any $t > t_0$,
% \begin{align*}
% \mathbb{P}\{ u(t,\cdot) \in \mathscr{K}(\Lambda) \} & \ge (1-
% \epsilon)\mathbb{P}\{ \Norm{u(t-t_0,\cdot)}_{\hat\rho} < \Lambda
% \}.
% \end{align*}
% We now apply Chebyshev's inequality and \eqref{E:InitData} to see that for all $t > t_0$ that there exists a $C>0$ such that
% \[
% \mathbb{P}\{ \Norm{u(t-t_0,\cdot)}_{\hat\rho} \ge \Lambda
% \} \le \frac{\mathbb{E}\left( \Norm{u(t-t_0,\cdot)}_{\hat\rho}^2 \right)}{\Lambda^2} \le C\frac{\mathcal{G}_{\hat{\rho}, q}\left( t - t_0 \right) }{\Lambda^2} \le \frac{C}{\Lambda^2}.
% \]
% Thus, by enlarging $\Lambda$ if necessary,
% \[
% \mathbb{P}\{ \Norm{u(t-t_0,\cdot)}_{\hat\rho} \ge \Lambda
% \}  \le \epsilon, \quad \forall \; t \ge t_0,
% \]
% thus completing the proof of \eqref{E:cpt}.
% % Since we assumed that $u$ is bounded in probability in $L^2_{\hat \rho}$ then for and $\epsilon >0$ we can fix $R>0$ such that $\mathbb{P}\{ |u(t,\cdot)|_{\hat\rho} \ge R \} < \epsilon$ and then taking $\Lambda > R$ such that $ \dfrac{\sin^q(\alpha \pi)}{\pi^q \Lambda^q}\bigg( 1+ R^q \bigg) < \epsilon$ concludes the claim.
%
% %\newpage
% % In this proof we will use $\stackrel{\mathcal{L}}{=}$ to denote equality in law, otherwise, the equality holds almost surely. Let $t_0>0$ and consider the following restarted equation:
% % \begin{equation} \label{E:SHErandint}
% %  \begin{cases}
% %   \dfrac{\partial v}{\partial t}(t,x)-\dfrac{1}{2}\Delta v(t,x) = b(x,v(t,x))\dot W_{t_0/2}(t,x) & \text{$x\in \R^d$ , $t>0$}, \\
% %   v(0,\cdot) = u(t_0/2,\cdot),                                                               &
% % \end{cases}
% % \end{equation}
% %where $\dot W_{t_0/2}(t,x) = \dot W(t + t_0/2,x)$ denotes the time shifted noise. \textcolor{blue}{Should the time shifted noise be equality in law?} We note that \eqref{E:uInLrho} implies that this random initial condition, $u(t_0/2,\cdot)$ satisfies \eqref{A:icon} and thus the solution to \eqref{E:SHErandint} exists and we denote the solution as $v$. By the Markov property of the solution to \eqref{E:SHE}, we may say that for any $t \ge 0$,  $u(t + t_0/2,x) \stackrel{\mathcal{L}}{=} v(t,x)$, where $\mathcal{L}$ means the equality is in law. In other words, for all $t \ge 0$, the following holds:
% % \begin{align*}
% %  u(t + t_0/2,x ) \stackrel{\mathcal{L}}{=} v(t,x) = (G(t,\cdot) * u(t_0/2,\cdot))(x) + \int_0^{t}\int_{\R^d} G(t - s, x-y)b(y,v(s, y)) W_{t_0/2}(\ud s, \ud y).
% % \end{align*}
% %We will use the notation $Y_v$ and $\mathcal{G}_{\hat\rho, q}^v$ where
% %  \[
% %    Y_v(s,y):= \int_0^s \int_{\R^d} (s-r)^{-\alpha}G(s-r,y-z)b(z,v(r,z))W_{t_0/2}(\ud r,\ud z),
% %  \]
% % and
% % \begin{align*}
% %  \mathcal{G}_{\hat\rho, q}^v(s) &= \int_{\R^d} \left(\int_{\R^d} G(s,x-y)\Norm{v(0,y)}_q\ud y\right)^2 \hat\rho(x) \ud x \\
% %  &= \int_{\R^d} \left(\int_{\R^d} G(s,x-y)\Norm{u(t_0/2,y)}_q\ud y\right)^2 \hat\rho(x) \ud x.
% % \end{align*}
% % We define $\mathcal{G}_{\hat\rho, q}^u(s)$ and $Y_u$ in an equivalent manner. We wish to be able to apply the factorization result \eqref{E:1fac} and also the bound in \eqref{E:yineq} for $Y_v$. This requires us to show that $v(0,\cdot) = u(t_0/2, \cdot)$ satisfies \eqref{E:G} and \eqref{E:GH-int'} with $\hat{\rho}$.  Indeed, H\"older's inequality and the admissible weight property imply that
% % \begin{align*}
% %  \mathcal{G}_{\hat\rho, q}^v(s)  &= \int_{\R^d} \left(\int_{\R^d} G(s,x-y)\Norm{u(t_0/2,y)}_q\ud y\right)^2 \hat\rho(x) \ud x \\
% %  &\le \int_{\R^d} \int_{\R^d} G(s,x-y)\Norm{u(t_0/2,y)}_q^2\ud y  \; \hat\rho(x) \ud x \\
% %  & \le C(t_0) \int_{\R^d} \Norm{u(t_0/2,y)}_q^2 \hat\rho(y) \ud y.
% % \end{align*}
% %Now by \eqref{E:mb0},
% % \begin{align*}
% %  C(t_0) \int_{\R^d} \Norm{u(t_0/2,y)}_q^2 \hat\rho(y) \ud y &\le C(t_0) \int_{\R^d} \Big(G(t_0/2,\cdot) * (1+\Norm{\mu}_q)\Big)^2(y) \hat\rho(y) \ud y \\
% %  &= C(t_0) \mathcal{G}_{\hat\rho, q}^u(t_0/2) < \infty,
% % \end{align*}
% %where the finiteness follows from \eqref{E:InitData}. Finally, this and the continuity of $\mathcal{H}_\alpha(s)$ for $s \in [0,t_0/2]$ shows that,
% % \[
% %  \int_0^{t_0/2} \left(\mathcal{G}_{\hat\rho, q}^v(s) \mathcal{H}_{\alpha}(s)\right)^{q/2} \ud s \le \mathcal{G}_{\hat\rho, q}^u(t_0/2)^{q/2} \int_0^{t_0/2} \left( \mathcal{H}_{\alpha}(s)\right)^{q/2} \ud s < \infty.
% % \]
% %Thus we are justified in saying that by \eqref{E:1fac},
% %\begin{equation}
% %u(t_0,x) \stackrel{\mathcal{L}}{=} v(t_0/2,x) = \Big(G(t_0/2,\cdot) * u(t_0/2,\cdot)\Big)(x) +\dfrac{\sin(\alpha \pi)}{\pi} \left[F_\alpha Y_v\right](t_0/2,x).
% %\end{equation}
% %We now define the set $\mathscr{K} (\Lambda)$ as
% % \[
% % \mathscr{K} (\Lambda) = \{ \big(G(t_0,\cdot)*y(\cdot)\big)(x)+(F_\alpha h)(x),
% % \quad \text{where $\Norm{y}_{\hat\rho} \le \Lambda$  and  $\Norm{h}_{L^q(0,1,L^2_{\hat \rho})} \le \Lambda $} \}.
% % \]
% %Then by Proposition \ref{P:F-Comp}, $\mathscr{K} (\Lambda)$ is relatively compact in $L^2_{ \rho}$ for $\Lambda >0$. Moreover, by \eqref{E:yineq} and Chebyshev's inequality we have
% % \begin{align}
% % \mathbb{P}\{ u(t_0,\cdot) \not \in \mathscr{K}(\Lambda) \}
% % &= \mathbb{P}\{ v(t_0/2,\cdot) \not \in \mathscr{K}(\Lambda) \} \nonumber \\
% % & \le \mathbb{P}\left[ \left(\int_{0}^{t_0/2} \Norm{Y_v(s,\cdot)}_{\widehat{\rho}}^q \: \ud s\right)^{1/q} > \dfrac{\pi \Lambda}{\sin(\alpha \pi )} \right] + \mathbb{P}\left[ \Norm{u(t_0/2,\cdot)}_{\widehat{\rho}} > \Lambda \right] \nonumber \\
% % & \le \dfrac{\sin^q(\alpha \pi)}{\pi^q \Lambda^q}\E \int_{0}^{t_0/2} \Norm{Y_v(s,\cdot)}_{\hat \rho}^q \: \ud s + \frac{1}{\Lambda^2} \mathbb{E}\left(\Norm{u(t_0/2,\cdot}_{\widehat{\rho}}^2\right) \nonumber \\
% % & \le \dfrac{\sin^q(\alpha \pi)}{\pi^q \Lambda^q} \Theta\int_0^{t_0/2} \left(\mathcal{G}_{\hat\rho, q}^v(s) \mathcal{H}_{\alpha}(s)\right)^{q/2} \ud s + \frac{1}{\Lambda^2} \mathbb{E}\left(\Norm{u(t_0/2,\cdot}_{\widehat{\rho}}^2\right) < \infty \label{E:Restartbdd}
% % \end{align}
% %With all this said, \eqref{E:Restartbdd} implies that
% % \begin{align*}
% % \mathbb{P}\{ u(t_0,\cdot) \in \mathscr{K}(\Lambda) \}
% % & \ge 1- \dfrac{\sin^q(\alpha \pi)}{\pi^q \Lambda^q} \left(\Theta\int_0^{t_0/2} \left(\mathcal{G}_{\hat\rho, q}^v(s) \mathcal{H}_{\alpha}(s)\right)^{q/2} \ud s\right) - \frac{1}{\Lambda^2} \mathbb{E}\left(\Norm{u(t_0/2,\cdot)}_{\rho}^2\right)
% % \end{align*}
% % Therefore, we can choose $\Lambda > 0$ big enough such that
% %  \[
% %   \dfrac{\sin^q(\alpha \pi)}{\pi^q \Lambda^q} \left(\Theta\int_0^{t_0/2} \left(\mathcal{G}_{\hat\rho, q}^v(s) \mathcal{H}_{\alpha}(s)\right)^{q/2} \ud s\right) + \frac{1}{\Lambda^2} \mathbb{E}\left(\Norm{u(t_0/2,\cdot)}_{\rho}^2\right) < \epsilon,
% %  \]
% % then we get that
% % \begin{align*}
% % \mathbb{P}\{ u(t_0,\cdot) \in \mathscr{K}(\Lambda) \} & \ge 1- \epsilon.
% % \end{align*}
% % We can now deduce from this and the semigroup property of the solution to the heat equation that
% % \begin{align*}
% % \mathbb{P}\{ u(t,\cdot) \in \mathscr{K}(\Lambda) \} & \ge (1-
% % \epsilon)\mathbb{P}\{ \Norm{u(t-t_0/2,\cdot)}_{\hat\rho} < \Lambda
% % \}.
% % \end{align*}
% % We now apply Chebyshev's inequality and \eqref{E:InitData} to see that for all $t > t_0/2$ that there exists a $C>0$ such that
% %  \[
% %   \mathbb{P}\{ \Norm{u(t-t_0/2,\cdot)}_{\hat\rho} \ge \Lambda
% %   \} \le \frac{\mathbb{E}\left( \Norm{u(t-t_0/2,\cdot)}_{\hat\rho}^2 \right)}{\Lambda^2} \le C\frac{\mathcal{G}_{\hat{\rho}, q}\left( t - \frac{t_0}{2} \right) }{\Lambda^2} \le \frac{C}{\Lambda^2}.
% %  \]
% % Thus, by enlarging $\Lambda$ if necessary,
% %  \[
% %   \mathbb{P}\{ \Norm{u(t-t_0/2,\cdot)}_{\hat\rho} \ge \Lambda
% %   \}  \le \epsilon, \quad \forall \; t \ge t_0/2,
% %  \]
% %thus completing the proof of \eqref{E:cpt}.
% %% Since we assumed that $u$ is bounded in probability in $L^2_{\hat \rho}$ then for and $\epsilon >0$ we can fix $R>0$ such that $\mathbb{P}\{ |u(t,\cdot)|_{\hat\rho} \ge R \} < \epsilon$ and then taking $\Lambda > R$ such that $ \dfrac{\sin^q(\alpha \pi)}{\pi^q \Lambda^q}\bigg( 1+ R^q \bigg) < \epsilon$ concludes the claim.
%\end{proof}
%%%%%%%%%%%%%%%%%%%%%%%%%%%%%%%%%%%%%%
% }}}

% {{{ All references: New ones with biber
\addcontentsline{toc}{section}{References}
\printbibliography[title={References}]
% }}}
\end{document}
\section{References}
% \cite{assing.manthey:03:invariant}
% \cite{carmona.molchanov:94:parabolic}
% \cite{cerrai:01:second}
% \cite{chen.dalang:15:moments}
% \cite{chen.hu.ea:21:regularity}
% \cite{chen.huang:19:comparison}
% \cite{chen.kim:19:nonlinear}
% \cite{dalang.khoshnevisan.ea:09:minicourse}
% \cite{dalang.quer-sardanyons:11:stochastic}
% \cite{da-prato.zabczyk:92:stochastic}
% \cite{da-prato.g-atarek.ea:92:invariant}
% \cite{tessitore.zabczyk:98:invariant}
% \cite{da-prato.zabczyk:96:ergodicity}
% \cite{walsh:86:introduction}
% \cite{olver.lozier.ea:10:nist}
% \cite{misiats.stanzhytskyi.ea:20:invariant}
% \cite{misiats.stanzhytskyi.ea:16:existence}
% \cite{maslowski.seidler:99:on}
% \cite{grafakos:14:modern}
